% concatenated from arxiv.tex
\documentclass[a4paper,11pt]{article}
\usepackage[T1]{fontenc}
\usepackage{lmodern}
\usepackage{amsmath,amssymb,amsthm}
\usepackage{hyperref}
\usepackage[capitalise,nameinlink,noabbrev]{cleveref}
\crefname{assumption}{Assumption}{Assumptions}
\crefname{conjecture}{Conjecture}{Conjectures}
\crefname{lemma}{Lemma}{Lemma}
\crefname{proposition}{Proposition}{Propositions}
\crefname{equation}{}{}
\crefname{algocf}{Algorithm}{Algorithms}
\crefname{Question}{Question}{Questions}
\newtheorem{example}{Example}
\newtheorem{remark}{Remark}
\newtheorem{question}{Question}
\newtheorem{assumption}{Assumption}
\newtheorem{theorem}{Theorem}
\newtheorem{lemma}{Lemma}
\newtheorem{definition}{Definition}
\usepackage[a4paper, margin=3cm]{geometry}

\usepackage{graphicx}
\usepackage{amsfonts,cite,float,url,booktabs}
\usepackage[section]{placeins}
\usepackage[algo2e,ruled,linesnumbered]{algorithm2e}
\newcommand{\argmin}{\mathop{\arg\min}}
\newcommand{\Ttran}{\mathsf{T}}
\newcommand{\Htran}{\mathsf{H}}
\newcommand{\de}{\,\mathrm{d}}
\newcommand{\mi}{\!\cdot\!\mathsf{i}}
\newcommand{\Rq}{\mathsf{Rf}}
\newcommand{\defi}{:=}
\newcommand{\order}{\mathcal{O}}
\newcommand{\R}{\mathbb{R}}
\newcommand{\C}{\mathbb{C}}
\newcommand{\rk}{\mathrm{rank}}
\newcommand{\region}{\mathcal{D}}
\newcommand{\range}{\mathsf{range}}
\newcommand{\fig}{eps}
\newcommand{\Crval}{C_{\mathsf{val}}}
\newcommand{\Crvec}{C_{\mathsf{vec}}}
\newcommand{\kval}{\kappa_{\mathsf{val}}}
\newcommand{\kvec}{\kappa_{\mathsf{vec}}}
\newcommand{\figsizeD}{0.4\textwidth}
\providecommand{\spa}[1]{\mathsf{span}\{#1\}}
\providecommand{\abs}[1]{\lvert#1\rvert}
\providecommand{\norm}[1]{\lVert#1\rVert}
\providecommand{\Bigabs}[1]{\Bigl\lvert#1\Bigr\rvert}
\providecommand{\bignorm}[1]{\bigl\lVert#1\bigr\rVert}
\providecommand{\Bignorm}[1]{\Bigl\lVert#1\Bigr\rVert}


\begin{document}
\title{Stabilizing the Rayleigh--Ritz procedure by randomization}
\author{ 
    Nian Shao\thanks{
         Institute of Mathematics, EPF Lausanne, 1015 Lausanne, Switzerland
         (\href{mailto:nian.shao@epfl.ch}{nian.shao@epfl.ch}).
     }
}
\maketitle

\begin{abstract}
Extracting approximate eigenpairs from a prescribed subspace is of fundamental importance in eigenvalue computation.
While projecting the target eigenvector onto the subspace yields satisfactory accuracy, extracting an approximate eigenpair that attains a comparable convergence rate has remained a long-standing open problem.
Although the standard Rayleigh--Ritz procedure is widely used for this purpose, it may suffer from deteriorated convergence of Ritz values and may even fail to produce convergent Ritz vectors.
In this paper, we address this long-standing open problem by introducing a randomized Rayleigh--Ritz procedure whose output converges at a rate similar to the ideal projection.
Our analysis requires only the simplicity of the target eigenvalue and extends naturally to nonlinear eigenvalue problems.


\end{abstract}

\medskip\textbf{Keywords:} Rayleigh--Ritz procedure, randomized algorithms,  eigenvalue problems

\medskip\textbf{AMS subject classifications (2020):} 65F15, 65H17, 68W20

\section{Introduction}
This paper is concerned with the following eigenpair extraction problem.
\begin{question}
    \label{Q}
    Let $(\lambda,v)$ be a simple eigenpair of $A_{0}\in\C^{n\times n}$. Given an $m$-dimensional subspace $\mathcal{W}_{\epsilon}$, where $1<m\ll n$ and 
\begin{equation*}
    \epsilon\defi \angle(v,\mathcal{W}_{\epsilon})
\end{equation*}
denotes the canonical angle between $v$ and $\mathcal{W}_{\epsilon}$. How can we reliably extract a good approximation of the eigenpair $(\lambda,v)$ from the prescribed subspace $\mathcal{W}_{\epsilon}$? 
\end{question}

\cref{Q} is of fundamental importance in eigenvalue computation, since such subspaces may arise from a variety of methods, such as Krylov subspace methods and reduced-basis techniques.
Unlike linear systems, where the optimal solution can be computed directly via a least squares formulation, extracting an eigenpair from a given subspace is considerably more challenging.
While the (standard) Rayleigh--Ritz procedure \cite[Chap.~4.4]{Stewart2001} is widely used for this purpose, Stewart \cite[p.~282]{Stewart2001} observes that ``we have no guarantee that the result approximates the desired eigenpair, \ldots, this difficulty has no easy answer'', highlighting \cref{Q} as a long-standing challenge.
Indeed, unless $A_{0}$ is Hermitian and $\lambda$ is an exterior eigenvalue, several stability issues may arise in the standard Rayleigh--Ritz procedure, including deterioration in the convergence of Ritz values and possible failure of convergence of Ritz vectors; see \Cref{sec:stab} for details.


In this paper, we introduce a \emph{randomized Rayleigh--Ritz} (RRR) procedure (\cref{alg:RRR-linear})  that provides an answer to \cref{Q}.
Before describing the RRR procedure, we briefly recall the standard Rayleigh--Ritz procedure. For the eigenvalue problem associated with $A_{0}$, it computes Ritz pairs $(\mu_{\epsilon}, W_{\epsilon} y_{\epsilon})$  by solving a compressed eigenvalue problem
\begin{equation*}
    W_{\epsilon}^{\Htran}A_{0}W_{\epsilon}y_{\epsilon}=\mu_{\epsilon} y_{\epsilon},
\end{equation*}
where $W_{\epsilon}$ is an orthonormal basis of $\mathcal{W}_{\epsilon}$.
In the RRR procedure, instead of imposing a Galerkin condition, we impose a Petrov--Galerkin condition using a complex Gaussian random matrix 
$\Omega\in\C^{n\times m}$ and solve the following compressed generalized eigenvalue problem:
\begin{equation*}
    \Omega^{\Htran}A_{0}W_{\epsilon}y_{\epsilon}=\mu_{\epsilon}\Omega^{\Htran}W_{\epsilon} y_{\epsilon}.
\end{equation*}
To distinguish these quantities from standard Ritz values and vectors, we refer to $\mu_{\epsilon}$ and $W_{\epsilon}y_{\epsilon}$ satisfying this compressed generalized eigenvalue problem as \emph{randomized Ritz values} and \emph{randomized Ritz vectors}, respectively. We then refine randomized Ritz values using the usual Rayleigh quotient
\begin{equation*}
    \rho_{\epsilon} = \frac{(W_{\epsilon}y_{\epsilon})^{\Htran}A_{0}(W_{\epsilon}y_{\epsilon})}{(W_{\epsilon}y_{\epsilon})^{\Htran}(W_{\epsilon}y_{\epsilon})}.
\end{equation*}
Our theoretical analysis shows that, for a \emph{simple} eigenpair $(\lambda,v)$ of $A_{0}$, with high probability, 
\begin{equation*}
    \abs{\lambda-\rho_{\epsilon}} = \order(\epsilon^{\gamma}),\quad 
    \abs{\lambda-\mu_{\epsilon}} = \order(\epsilon)
    \quad \text{and} \quad
    \angle(v, W_{\epsilon}y_{\epsilon}) = \order(\epsilon),
\end{equation*}
where $\gamma=2$ for a Hermitian $A_{0}$ and $\gamma=1$ otherwise.
We emphasize that the $\order(\cdot)$ notation is employed solely for notational convenience; all probabilistic bounds in this paper are fully non-asymptotic.
In particular, if there exists a family of subspaces $\mathcal{W}_{\epsilon}$ satisfying $\epsilon\to 0$, then up to a constant independent of the subspaces $\mathcal{W}_{\epsilon}$, the convergence rates of randomized Ritz vectors and their corresponding refined randomized Ritz values match that of the ideal projections $W_{\epsilon}W_{\epsilon}^{\Htran}v$ and their Rayleigh quotients, respectively.
To the best of our knowledge, no existing method is known to achieve this.






\begin{algorithm2e}[H]
    \caption{RRR procedure for standard eigenvalue problems}
    \label{alg:RRR-linear}
    \KwIn{$A_{0}\in\C^{n\times n}$ and an orthonormal basis $W_{\epsilon}\in\C^{n\times m}$\;}
    Draw a complex Gaussian random matrix $\Omega\in\C^{n\times m}$\; 
    Solve the compressed generalized eigenvalue problem of the pencil $\Omega^{\Htran}A_{0}W_{\epsilon}-\xi \Omega^{\Htran}W_{\epsilon}$ and select wanted randomized Ritz pairs $(\mu_{\epsilon},W_{\epsilon}y_{\epsilon})$\;
    Refine randomized Ritz values by $\rho_{\epsilon} = \frac{(W_{\epsilon}y_{\epsilon})^{\Htran}A_{0}(W_{\epsilon}y_{\epsilon})}{(W_{\epsilon}y_{\epsilon})^{\Htran}(W_{\epsilon}y_{\epsilon})}$\;
    \Return{Randomized Ritz pairs $(\mu_{\epsilon},W_{\epsilon}y_{\epsilon})$ and refined randomized Ritz values $\rho_{\epsilon}$\;}
\end{algorithm2e}

We remark that \cref{alg:RRR-linear} is closely related to the sketched Rayleigh--Ritz procedure proposed in \cite[Sec.~6]{Nakatsukasa2024}. However, there are several differences between the two methods. In the sketched Rayleigh--Ritz procedure, a random matrix is used as an oblivious subspace embedding, and oversampling is required, whereas the RRR procedure does not require oversampling. Moreover, the sketched Rayleigh--Ritz procedure provides only \emph{a posteriori} estimates, while we establish \emph{a priori} convergence results.
Other related algorithms for extracting eigenpairs include the refined \cite{Jia1997} and harmonic \cite{Morgan1991} Rayleigh--Ritz procedures. Compared with the standard Rayleigh--Ritz procedure, these variants require additional computations, such as SVD or QR decomposition of $n\times m$ matrices. From a theoretical perspective, they can be interpreted as performing an alternating optimization over $\mu$ and $y$, which work well empirically for standard eigenvalue problems. However, as we will show in \Cref{sec:stab}, both Ritz values and Ritz vectors may fail to converge for generalized (and nonlinear) eigenvalue problems. Consequently, rigorous convergence guarantees for the refined and harmonic Rayleigh--Ritz procedures are currently unavailable.



The remainder of this paper is organized as follows.
In \Cref{sec:stab}, we present several examples illustrating the stability issues of the standard Rayleigh--Ritz procedure.
In \Cref{sec:algo}, we formally describe the RRR procedure for nonlinear eigenvalue problems.
\Cref{sec:theory} contains our main theoretical results, including the linear convergence of randomized Ritz pairs (\cref{thm:ritzpair}) and the linear or quadratic convergence of refined randomized Ritz values (\cref{thm:Rf,thm:nls}).
As in the standard eigenvalue problem setting, our probabilistic non-asymptotic convergence results require only that $\lambda$ is a simple eigenvalue of the \emph{original} problem.
Numerical experiments are presented in \Cref{sec:numexp}.


\paragraph{Notation}
Throughout this paper, $\norm{\cdot}$ denotes the Euclidean norm for vectors and the spectral norm for matrices.
We use $e_{i}$ to denote the $i$-th canonical basis vector.
A matrix $\Omega$ is called a complex Gaussian random matrix if the real and imaginary parts of its entries are independent and identically distributed according to $\mathcal{N}(0,1/2)$.
The canonical angle between $\spa{v}$ and $\range(W)$ is denoted by $\angle(v,W)$ or $\angle(v,\range(W))$.
For a subspace $\mathcal{W}\subset\C^{n}$, we denote its orthogonal complement by $\mathcal{W}^{\perp}$, so that $\mathcal{W}\oplus \mathcal{W}^{\perp} = \C^{n}$.
A matrix $Q\in\C^{n\times m}$ is called orthonormal if its columns form an orthonormal basis, that is, $Q^{\Htran}Q = I_{m}$, where $I_{m}$ is the $m\times m$ identity matrix.
For a matrix-valued function $A(\xi)$ and a nonempty bounded domain $\region$, we define $
\norm{A}_{\region} := \sup_{\xi\in\region}\norm{A(\xi)}$.




\section{Stability issues of standard Rayleigh--Ritz procedure}
\label{sec:stab}
In this section, we present several examples illustrating the stability issues of the standard Rayleigh--Ritz procedure. Throughout this section, as well as in all results of this paper, we assume that $(\lambda, v)$ is a simple eigenpair.

We begin with a Hermitian matrix $A_{0}$. When $\lambda$ is the largest (or smallest) eigenvalue, the minimax principle ensures that the second eigenvalue of the compressed matrix $W_{\epsilon}^{\Htran} A_{0} W_{\epsilon}$ is bounded by the second eigenvalue of $A_{0}$. If $\mu_{\epsilon}$ converges to $\lambda$, then it remains well separated from the remaining eigenvalues of $W_{\epsilon}^{\Htran} A_{0} W_{\epsilon}$. Consequently, Ritz values converge quadratically, while Ritz vectors converge linearly; see, for instance, \cite[Sec.~4.3.2]{Saad2011}.

In contrast, for an interior eigenvalue of a Hermitian matrix, the convergence of Ritz values may deteriorate to linear, and Ritz vectors may even fail to converge. This occurs because the target Ritz value can become arbitrarily close to other eigenvalues of $W_{\epsilon}^{\Htran} A_{0} W_{\epsilon}$.
\begin{example}[Interior eigenvalue of a Hermitian matrix: linear convergence of Ritz values and failure of convergence of Ritz vectors]
    \label{eg:H}
    Consider the eigenvalue $0$ of the matrix  
    \begin{equation*}
        A_{0} = \begin{bmatrix}
            -1 & 0 & 0 \\ 
            0 & 0 & 0\\ 
            0 & 0 & 1
        \end{bmatrix}
        \quad\text{and}\quad 
        W_{\epsilon} = \begin{bmatrix}
            \epsilon/\sqrt{2} & 1/\sqrt{2}\\ 
            \sqrt{1-\epsilon^{2}} & 0 \\ 
            \epsilon/\sqrt{2} & -1/\sqrt{2}
        \end{bmatrix}.
    \end{equation*}
    Then the compressed matrix is 
    \begin{equation*}
        W_{\epsilon}^{\Htran}A_{0}W_{\epsilon} = \begin{bmatrix}
            0 & -\epsilon\\ 
            -\epsilon & 0
        \end{bmatrix}.
    \end{equation*}
    Ritz values are $\pm\epsilon$, yielding only linear convergence.
    Corresponding Ritz vectors are $W_{\epsilon}[1,-1]^{\Ttran}$ and $W_{\epsilon}[1,1]^{\Ttran}$, which provide no accuracy for target eigenvectors.
\end{example}


For a non-Hermitian matrix $A_{0}$, as shown in \cite{Jia2001}, applying Elsner's theorem 
to the compressed matrix ensures the convergence of Ritz values.  However, the convergence rate may deteriorate substantially to $\order(\epsilon^{1/m})$, since the compressed matrix can become (nearly) defective.
\begin{example}[Non-Hermitian matrix: $\order(\epsilon^{1/m})$ convergence on Ritz values]
    \label{eg:nH}
    Consider the eigenvalue $0$ of the matrix 
\begin{equation*}
    A_{0} = \begin{bmatrix}
        0 & 1 & 0\\ 
        0 & 0 & I_{m-1}\\ 
        0 & 1 & 0
    \end{bmatrix}
    \quad\text{and}\quad 
    W_{\epsilon}=\begin{bmatrix}
        \sqrt{1-\epsilon^{2}} & 0\\ 
        0 & I_{m-1} \\ 
        \epsilon & 0
    \end{bmatrix}.
\end{equation*}
The compressed matrix is 
\begin{equation*}
    W_{\epsilon}^{\Htran}A_{0}W_{\epsilon}
     = 		\begin{bmatrix}
			0        & \sqrt{1-\epsilon^2} + \epsilon & 0       \\
			0        & 0                              & I_{m-2} \\
			\epsilon & 0                              & 0
		\end{bmatrix},
\end{equation*}
whose eigenvalues are approximately of magnitude $\order(\epsilon^{1/m})$ as $\epsilon\to 0$.
\end{example}


An even more severe situation arises for generalized eigenvalue problems, where Ritz values may fail to converge. This occurs because the compressed matrix pencil can be $\order(\epsilon)$-close to singular even if the original pencil is regular. Consequently, the convergence result for standard eigenvalue problems cannot be extended to the generalized case. 

\begin{example}[Generalized eigenvalue problem: Ritz values fail to converge\footnote{Unfortunately, this example shows that the claimed unconditional convergence results of Ritz values and refined Ritz vectors in \cite{Jia2025} does not hold in general.}]
    \label{eg:G}
    Consider the matrix pencil $A_{0}-\xi A_{1}$, where
\begin{equation*}
    A_{0}= \begin{bmatrix}
        0 & 0 & 1\\
        0 & 2 & 0\\
        0 & 3 & 0
    \end{bmatrix}
    \quad\text{and}\quad
    A_{1}= \begin{bmatrix}
        0 & 0 & 1\\
        1 & 0 & 0\\
        0 & 1 & 0
    \end{bmatrix}.
\end{equation*}
Since $\det (A_{0}-\xi A_{1}) = -\xi(1-\xi)(3-\xi)$, the pencil is regular with three simple eigenvalues $0$, $1$ and $3$, associated with eigenvectors $[1,0,0]^{\Ttran}$, $[0,0,1]^{\Ttran}$ and $[2,3,0]^{\Ttran}$, respectively.
Consider the following trail subspace and its corresponding compressed pencil
\begin{equation*}
    W_{\epsilon} = \begin{bmatrix}
        \sqrt{1-\epsilon^{2}} & 0\\
        0 & 1\\
        \epsilon & 0
    \end{bmatrix}
    \quad\text{and}\quad
    W_{\epsilon}^{\Htran}(A_{0}-\xi A_{1})W_{\epsilon} = \begin{bmatrix}
        \epsilon &0\\
        0& 1
    \end{bmatrix}
    \begin{bmatrix}
        1-\xi & 3-\xi\\
        -\xi & 2
    \end{bmatrix}\begin{bmatrix}
        \sqrt{1-\epsilon^{2}} & 0\\
        0& 1
    \end{bmatrix}.
\end{equation*}
Ritz values are $2$ and $-1$ for every $\epsilon>0$, which fail to converge to an eigenvalue of $A_{0}-\xi A_{1}$ as $\epsilon \to 0$.
\end{example}

We summarize the instabilities of the standard Rayleigh--Ritz procedure and compare them with those of the RRR procedure in \cref{tab:cmp}.
The corresponding convergence results for the latter will be established in \cref{thm:nls,thm:Rf,thm:ritzpair}.


\begin{table}[H]
    \centering
    \caption{Convergence guarantees of the standard and randomized Rayleigh--Ritz procedures with $\epsilon = \angle(v, W_{\epsilon})$. 
A cross `$\times$' indicates that the convergence may fail, while the question mark `?' denotes that the result is currently unclear. 
An eigenvalue problem $A(\cdot)$ is called Hermitian (at $\lambda$) if the matrix $A(\lambda)$ is Hermitian for the target eigenvalue $\lambda$. 
All convergence statements for the RRR procedure are probabilistic and non-asymptotic.}
    \begin{tabular}{lccccc}
        \toprule
        & \multicolumn{2}{c}{Rayleigh--Ritz}  & \multicolumn{3}{c}{randomized Rayleigh--Ritz} \\ \cmidrule(lr){2-3} \cmidrule(lr){4-6}
        & value & vector  & refined value & value & vector\\ \midrule
        Linear Hermitian (exterior) & $\order(\epsilon^{2})$ & $\order(\epsilon)$ & $\order(\epsilon^{2})$  & $\order(\epsilon)$ & $\order(\epsilon)$\\ 
        Linear Hermitian (interior) & $\order(\epsilon)$ & $\times$ & $\order(\epsilon^{2})$  & $\order(\epsilon)$ & $\order(\epsilon)$\\ 
        Nonlinear Hermitian& ? & $\times$ & $\order(\epsilon^{2})$  & $\order(\epsilon)$ & $\order(\epsilon)$\\ 
        Standard non-Hermitian & $\order(\epsilon^{1/m})$ & $\times$ & $\order(\epsilon)$  & $\order(\epsilon)$ & $\order(\epsilon)$\\ 
        Generalized non-Hermitian & $\times$ & $\times$ & $\order(\epsilon)$  & $\order(\epsilon)$ & $\order(\epsilon)$\\ 
        Nonlinear non-Hermitian & $\times$ & $\times$ & $\order(\epsilon)$  & $\order(\epsilon)$ & $\order(\epsilon)$\\
         \bottomrule
        \end{tabular}
    \label{tab:cmp}
\end{table}


\begin{remark}
For a non-Hermitian problem with left and right eigenvectors $u$ and $v$, suppose that we are given two subspaces $\mathcal{W}_{\varepsilon,L}$ and $\mathcal{W}_{\epsilon,R}$ such that $\varepsilon=\angle(u,\mathcal{W}_{\varepsilon,L})$ and $\epsilon=\angle(v,\mathcal{W}_{\epsilon,R})$.
We can apply the RRR procedure separately to obtain left and right randomized Ritz vectors $W_{\varepsilon,L}y_{\varepsilon,L}$ and $W_{\epsilon,R}y_{\epsilon,R}$.
With an additional refinement analogous to the two-sided Rayleigh quotient \cite{Schwetlick2012}, the accuracy of refined randomized Ritz values can be improved to $\order(\varepsilon\epsilon)$.
\end{remark}


At the end of this section, we comment on the instabilities of the standard Rayleigh--Ritz procedure and give an intuitive explanation of how randomization can stabilize it.
Consider a matrix $A_{0}$ with a simple eigenpair $(\lambda,v)$. Let $W$ be an orthonormal matrix such that $v\in\range(W)$.
The simplicity of $\lambda$ implies that zero is a simple singular value of $(A_{0}-\lambda I)W$.
However, the compressed matrix $W^{\Htran}(A_{0}-\lambda I)W$ does \emph{not} generally inherit this property unless $A_{0}$ is Hermitian and $\lambda$ is an exterior eigenvalue.
In fact, the compression may degenerate to a zero matrix in the Hermitian case (see \cref{eg:H}) or form a Jordan block in the non-Hermitian case (see \cref{eg:nH}).
As a result, the convergence of Ritz values may deteriorate and Ritz vectors may fail to converge.
In contrast, compressing $(A_{0}-\lambda I)W$ with an oblivious (complex) Gaussian random matrix preserves the simplicity of the zero singular value almost surely.
Specifically, in \cref{lem:kappa}, we show that, with high probability, the eigenvalue and eigenvector condition numbers of the compressed problem (at $\lambda$) are bounded by moderate multiples of those of the original problem.
These properties enable us to establish a probabilistic convergence result for the RRR procedure.




\section{RRR procedure for nonlinear eigenvalue problems}
\label{sec:algo}
In this section, we extend the RRR procedure to nonlinear eigenvalue problems (\cref{alg:RRR}).
Readers interested solely in standard eigenvalue problems may skip this section, with the exception of the perturbation result in \cref{lem:pert}.


\begin{assumption}
    \label{asp:A}
    Throughout this paper, we assume that $A(\cdot)$ is a regular holomorphic matrix-valued function defined on a nonempty bounded domain $\region\subset\C$ that admits an analytic extension to $\partial\region$.
    Here, regular means that $\det A(\xi_{0})\neq 0$ for some $\xi_{0}\in\region$.
\end{assumption}

The nonlinear eigenvalue problem \cite{Guttel2017} consists of finding eigenpairs $(\lambda,v)$ such that
\begin{equation*}
    A(\lambda) v = 0,\quad \lambda\in\region\quad\text{and}\quad v\in\C^{n}\setminus\{0\}.
\end{equation*}
In particular, if $0$ is a simple eigenvalue of the matrix $A(\lambda)$, and $u^{\Htran}A^{\prime}(\lambda)v\neq 0$ for a left eigenvector $u$, then the eigenvalue $\lambda$ of $A(\cdot)$ is called simple \cite[Prop.~1]{Neumaier1985}.



Similar to standard eigenvalue problems, we impose a Petrov--Galerkin condition on $A(\cdot)W_{\epsilon}$ using a complex Gaussian random matrix $\Omega\in\C^{n\times m}$ and solve the following compressed nonlinear eigenvalue problem:
\begin{equation}
    \label{eq:RRR}
    B_{\epsilon}(\mu_{\epsilon})y_{\epsilon} = 0,\quad 
    \mu_{\epsilon}\in\region,\quad 
    y_{\epsilon}\in\C^{m}\setminus\{0\},\quad\text{where}\quad 
    B_{\epsilon}(\xi) \defi \Omega^{\Htran}A(\xi)W_{\epsilon}.
\end{equation}

To refine randomized Ritz values, we may use the \emph{Rayleigh functional} \cite{Schwetlick2012}.
\begin{definition}
    \label{def:Rf}
    Let $\lambda$ be a simple eigenvalue of $A(\cdot)$ associated with an eigenvector $v$. Given two open sets $\region_{\lambda}\subset \region$ and $\region_{v}\subset \C^{n}$ such that $\lambda\in\region_{\lambda}$ and $v\in\region_{v}$, a functional $\Rq\colon \region_{v}\mapsto \region_{\lambda}$ is called a Rayleigh functional of $A(\cdot)$ if
    \begin{equation*}
        \Rq(\alpha w)=\Rq(w),\quad 
        w^{\Htran}A\bigl(\Rq(w)\bigr)w= 0,\quad 
        w^{\Htran}A^{\prime}\bigl(\Rq(w)\bigr)w\neq 0
    \end{equation*}
    holds for all $ \alpha\in\C\setminus\{0\}$ and $w\in\region_{v}$.
\end{definition}

For a simple eigenpair $(\lambda,v)$, if $v^{\Htran}A^{\prime}(\lambda)v\neq 0$, then the Rayleigh functional exists uniquely \cite[Thm.~21]{Schwetlick2012}. In particular, if $A(\lambda)$ is Hermitian, the simplicity of $\lambda$ implies the existence of $\Rq$. For a linear eigenvalue problem $A(\xi)=A_{0}-\xi A_{1}$, if $v^{\Htran}A_{1}v\neq 0$, then the Rayleigh functional is the usual Rayleigh quotient
\begin{equation*}
    \Rq(w) = \frac{w^{\Htran}A_{0}w}{w^{\Htran}A_{1}w}.
\end{equation*} 


If a Rayleigh functional is available,  randomized Ritz values can be refined by $\rho_{\epsilon}=\Rq(W_{\epsilon}y_{\epsilon})$, where $W_{\epsilon}y_{\epsilon}$ is a randomized Ritz vector.  
However, for non-Hermitian problems, a Rayleigh functional may not exist because $v^{\Htran}A^{\prime}(\lambda)v$ could be zero, such as the matrix pencil in \cref{eg:G}. In this situation, we can still perform the refinement via finding a stationary point 
\begin{equation}
    \label{eq:nls}
    \rho_{\epsilon} \in \Bigl\{ \rho\in\C \Bigm\vert \frac{\de }{\de \bar{\rho}} \norm{A(\rho)W_{\epsilon}y_{\epsilon}}^{2} = 0  \Bigr\},
\end{equation}
where $\bar{\rho}$ is the complex conjugate of $\rho$ and 
\begin{equation*}
    \frac{\de }{\de \bar{\rho}} \norm{A(\rho)W_{\epsilon}y_{\epsilon}}^{2} = \bigl(A^{\prime}(\rho)W_{\epsilon}y_{\epsilon}\bigr)^{\Htran}\bigl(A(\rho)W_{\epsilon}y_{\epsilon}\bigr)
\end{equation*}
is the Wirtinger derivative.
When $A(\xi)=A_{0}-\xi A_{1}$, this reduces to 
\begin{equation*}
    \rho_{\epsilon} = \frac{(W_{\epsilon}y_{\epsilon})^{\Htran}A_{1}^{\Htran}A_{0}(W_{\epsilon}y_{\epsilon})}{(W_{\epsilon}y_{\epsilon})^{\Htran}A_{1}^{\Htran}A_{1}(W_{\epsilon}y_{\epsilon})}.
\end{equation*}
For a general $A(\cdot)$, finding a stationary point of \cref{eq:nls} is closely related to a nonlinear least squares problem, which can be solved by the Gauss--Newton method or the Levenberg--Marquardt method \cite[Chap.~10]{Nocedal2006}.
If the set in \cref{eq:nls} admits multiple stationary points, additional selection criteria are required to identify an appropriate $\rho_{\epsilon}$, for instance, choosing the one closest to a prescribed shift.

We remark that those refinements described above, including both Rayleigh functional and stationary point approaches, typically do \emph{not} require additional matrix-vector multiplications with $A(\cdot)$, since $A(\cdot)W_{\epsilon}$ has already been computed when forming the compressed eigenvalue problem.




\begin{algorithm2e}[H]
    \caption{RRR procedure for nonlinear eigenvalue problems}
    \label{alg:RRR}
    \KwIn{$A(\cdot)\in\C^{n\times n}$ and an orthonormal basis $W_{\epsilon}\in\C^{n\times m}$\;}
    Draw a complex Gaussian random matrix $\Omega\in\C^{n\times m}$\;
    Solve the compressed nonlinear eigenvalue problem of $B_{\epsilon}(\xi)=\Omega^{\Htran}A(\xi)W_{\epsilon}$ and select wanted randomized Ritz pairs $(\mu_{\epsilon}, W_{\epsilon}y_{\epsilon})$\;
    Refine randomized Ritz values to obtain $\rho_{\epsilon}$ by evaluating the Rayleigh functional $\Rq(W_{\epsilon}y_{\epsilon})$ or finding a suitable stationary point of \cref{eq:nls}\;
    \Return{Randomized Ritz pairs $(\mu_{\epsilon},W_{\epsilon}y_{\epsilon})$ and refined randomized Ritz values $\rho_{\epsilon}$\;}
\end{algorithm2e}

We conclude this section with the following first-order perturbation results.
\begin{theorem}
    \label{lem:pert}
    Given a matrix-valued function $B(\cdot)\colon \region\mapsto\C^{m\times m}$ satisfying \cref{asp:A}. Let $\lambda$ be a simple eigenvalue of $B(\cdot)$ associated with left and right eigenvectors $x$ and $y$.
    Let $Z_{\perp}$ and $Y_{\perp}$ be orthonormal bases of $\spa{B^{\prime}(\lambda)y}^{\perp}$ and $\spa{y}^{\perp}$, respectively. Then $Z_{\perp}^{\Htran}B(\lambda)Y_{\perp}$ is nonsingular.

    Consider a perturbation $B+\Delta B$, where $\Delta B(\cdot)\colon \region\mapsto \C^{m\times m}$ is a holomorphic matrix-valued function that admits an analytic extension to $\partial\region$.
    Let 
    \begin{equation*}
        M = \max\bigl\{ \norm{B}_{\region},\norm{B^{\prime}}_{\region},\norm{B^{\prime\prime}}_{\region}\bigr\}
        \quad\text{and}\quad 
        \epsilon = \max\bigl\{\norm{\Delta B}_{\region},\norm{\Delta B^{\prime}}_{\region}, \norm{\Delta B^{\prime\prime}}_{\region}\bigr\}.
    \end{equation*}
    When $\epsilon$ is small enough, there exists an eigenpair $(\lambda +\Delta \lambda,y+\Delta y)$ of $B+\Delta B$ such that
    \begin{equation*}
        \abs{\Delta \lambda} \leq \kval(B,\lambda)\cdot\norm{\Delta B(\lambda)} +C \epsilon^{2} 
        \quad\text{and}\quad  
        \frac{\norm{\Delta y}}{\norm{y}} \leq \kvec(B,\lambda)\cdot\norm{\Delta B(\lambda)}+C\epsilon^{2},
    \end{equation*}
    where
    \begin{equation*}
        \kval(B,\lambda) \defi \frac{\norm{x}\norm{y}}{\abs{x^{\Htran}B^{\prime}(\lambda)y}}
        \quad \text{and}\quad 
        \kvec(B,\lambda)\defi \bignorm{(Z_{\perp}^{\Htran}B(\lambda)Y_{\perp})^{-1}}.
    \end{equation*} 
    Here, we use $y^{\Htran}\Delta y= 0$ as the normalization condition for $y+\Delta y$, and $C$ satisfies
    \begin{equation*}
        C\leq \widetilde{C}M\kval(B,\lambda)\cdot\bigl(\kval(B,\lambda)+\kvec(B,\lambda)\bigr)^{2}
    \end{equation*}
    for some universal constant $\widetilde{C}$.
\end{theorem}



The proof of \cref{lem:pert} is provided in \cref{app:pfpert}. 
We refer to the quantities $\kval(B,\lambda)$ and $\kvec(B,\lambda)$ as the eigenvalue and eigenvector condition numbers, respectively.
These names are consistent with the standard eigenvalue problem setting \cite[pp.~48--50]{Stewart2001}.
Compared with existing first-order perturbation results, such as asymptotic and non-asymptotic bounds for linear eigenvalue problems in \cite{Stewart1990}, and asymptotic results for (homogeneous) polynomial eigenvalue problems in \cite{Dedieu2000}, our result is slightly different.
Strictly speaking, \cref{lem:pert} is a non-asymptotic perturbation bound.
However, unlike~\cite{Stewart1990}, where the non-asymptotic control applies to the first-order term, our explicit bound controls the second-order remainder.
This formulation preserves the sharpness of the first-order expansion while simultaneously providing a non-asymptotic guarantee on higher-order terms, whose importance will be further explained in \cref{rmk:nonasymptotic}.




\section{Convergence analysis}
\label{sec:theory}
This section studies the convergence of \cref{alg:RRR}.
We remark that our results are also new for standard eigenvalue problems $A(\xi)=A_{0}-\xi I$.

 
\subsection{Exact recovery}
Let $(\lambda,v)$ be a simple eigenpair of $A(\cdot)$. Suppose that $v$ is contained in a subspace $\mathcal{W}$. In the following lemma, we show that this eigenpair can be recovered almost surely by solving the compressed nonlinear eigenvalue problem of 
\begin{equation}
    \label{eq:defB}
    B(\xi) \defi \Omega^{\Htran} A(\xi)W, 
    \quad \text{where}\quad  W=[v,W_{\perp}]
\end{equation}
forms an orthonormal basis of $\mathcal{W}$. 
\begin{lemma}
    \label{lem:exact}
    Let $(\lambda,v)$ be a simple eigenpair of $A(\cdot)$.  Given an arbitrary fixed full-rank matrix $W\in\C^{n\times m}$ satisfying $We_{1}=v$, the following statements hold for almost every $\Omega\in\C^{n\times m}$ with respect to Lebesgue measure:
    The matrix-valued function $B(\xi)=\Omega^{\Htran}A(\xi)W$ is regular, and $(\lambda,e_{1})$ is a simple eigenpair of $B(\cdot)$.
\end{lemma}
\begin{proof}
    Since $A(\cdot)$ is regular, we know that $A(\xi_{0})$ is nonsingular for some $\xi_{0}\in\region$. Then $A(\xi_{0})W$ is full rank, hence $\det B(\xi_{0}) = \det \bigl(\Omega^{\Htran}A(\xi_{0})W\bigr)\neq 0$ generically because this is a polynomial of the entries of $\Omega$.

Note that $B(\lambda)e_{1}=\Omega^{\Htran} A(\lambda)v=0$, hence $(\lambda,e_{1})$ is an eigenpair of $B(\cdot)$.
We now show its simplicity. Since $\lambda$ is a simple eigenvalue of $A(\cdot)$, the null space dimension of the matrix $A(\lambda)$ is one. It follows that, generically, the null space dimension of the matrix $B(\lambda)=\Omega^{\Htran}A(\lambda)W$ is also one, implying that $0$ is a simple eigenvalue of the matrix $B(\lambda)$.

Write $W=[v,W_{\perp}]$. Then $\rk(A(\lambda)W_{\perp})=m-1$. Let $u$ be a left eigenvector of $A(\cdot)$ corresponding to the eigenvalue $\lambda$. We have $u^{\Htran}A(\lambda)W_{\perp}=0$ and $u^{\Htran}A^{\prime}(\lambda)v\neq 0$. Hence, $A^{\prime}(\lambda)v\notin\range(A(\lambda)W_{\perp})$, which implies that $\rk[A^{\prime}(\lambda)v,A(\lambda)W_{\perp}]=m$. Consequently, for a generic $\Omega$, the matrix $\Omega^{\Htran}[A^{\prime}(\lambda)v,A(\lambda)W_{\perp}]\in\C^{m\times m}$ is nonsingular.
Let $x$ be a left eigenvector of $B(\cdot)$ corresponding to $\lambda$. Since $x^{\Htran}B(\lambda)=x^{\Htran}\Omega^{\Htran}A(\lambda)W=0$, we have
\begin{equation}
    \label{eq:pfexact}
x^{\Htran}\Omega^{\Htran}[A^{\prime}(\lambda)v,A(\lambda)W_{\perp}]
=[x^{\Htran}B^{\prime}(\lambda)e_{1},0_{1\times(m-1)}].
\end{equation}
Since $\Omega^{\Htran}[A^{\prime}(\lambda)v,A(\lambda)W_{\perp}]$ is nonsingular and $x\neq 0$, we conclude that $x^{\Htran}B^{\prime}(\lambda)e_{1}\neq 0$.
\end{proof}

The proof of \cref{lem:exact} provides useful intuition for the subsequent analysis of robust recovery.
Actually, we can directly derive an upper bound for the eigenvalue condition number $\kval(B,\lambda)$ from \cref{eq:pfexact}:
\begin{equation*}
    \kval(B,\lambda)
    = \frac{\norm{x}\norm{e_{1}}}{\abs{x^{\Htran}B^{\prime}(\lambda)e_{1}}}
    = \frac{\norm{x}}{\norm{x^{\Htran}\Omega^{\Htran}[A^{\prime}(\lambda)v,A(\lambda)W_{\perp}]}}
    \leq \bignorm{\bigl(\Omega^{\Htran}[A^{\prime}(\lambda)v,A(\lambda)W_{\perp}]\bigr)^{-1}}.
\end{equation*}
Although this bound is not as sharp as what we will establish in \cref{lem:kappa}, it suggests studying the matrix $[A^{\prime}(\lambda)v,A(\lambda)W_{\perp}]$ and exploiting the orthogonality between $\Omega x$ and $A(\lambda)W_{\perp}$.


\subsection{Probabilistic analysis on eigenvalue and eigenvector condition numbers}
Consider the ideal proxy $B(\cdot)$ defined in \cref{eq:defB}.
In this subsection, we derive probabilistic bounds for its eigenvalue condition number $\kval(B,\lambda)$ and eigenvector condition number $\kvec(B,\lambda)$, as introduced in \cref{lem:pert}. For this aim, we first recall some standard tail bounds for extreme singular values of complex Gaussian random matrices. These results can be found, for example, in \cite{Rudelson2010,Edelman1988}.
\begin{lemma}
    \label{lem:Gaussian}
    Let $\Omega\in\C^{n\times m}$ be a complex Gaussian random matrix. Then 
    \begin{equation*}
        \Pr\bigl(\norm{\Omega}\geq \sqrt{n}+\sqrt{m}+\sqrt{\ln(2/\delta)}\bigr)\leq \delta
        \quad\text{for all}\quad \delta\geq 0. 
    \end{equation*}
    When $m=n$, that is, $\Omega\in\C^{m\times m}$ is a square matrix (or a scalar), it holds that:
    \begin{equation*}
        \Pr\bigl(\norm{\Omega^{-1}}\geq \sqrt{m/\delta}\bigr)\leq \delta \quad\text{for all}\quad \delta \geq 0. 
    \end{equation*}
\end{lemma}

The following lemma shows that compressing the matrix $A(\lambda)W$ with a complex Gaussian random matrix $\Omega$ preserves both eigenvalue and eigenvector condition numbers with high probability.
We remark that replacing $\Omega$ with $W$, as in the standard Rayleigh--Ritz procedure, may completely destroy a well-conditioned eigenpair of $A(\cdot)$.

\begin{lemma}
    \label{lem:kappa}
    Let $u$ and $v$ be unit left and right eigenvectors of $A(\cdot)$ corresponding to a simple eigenvalue $\lambda$.
    Given an arbitrary fixed orthonormal matrix $W$ admitting $W=[v,W_{\perp}]\in\C^{n\times m}$, we define $B(\xi)=\Omega^{\Htran}A(\xi)W$, where $\Omega\in\C^{n\times m}$ is a complex Gaussian random matrix. Then 
    \begin{equation*}
        \begin{aligned}
                \kval(B,\lambda) &\leq  \frac{1}{\sqrt{\delta}}\cdot \kval(A,\lambda) &\text{holds with probability at least $1-\delta$,}\\ 
                \kvec(B,\lambda) &\leq  \frac{\sqrt{m-1}}{\sqrt{\delta}}\cdot \kvec(A,\lambda) &\text{holds with probability at least $1-\delta$.}\\ 
        \end{aligned}
    \end{equation*}
    Here, $\kval(A,\lambda)$ and $\kvec(A,\lambda)$ are eigenvalue and eigenvector condition numbers of $A(\cdot)$:
    \begin{equation}
        \label{eq:defkappaA}
        \kval(A,\lambda)\defi \frac{\norm{u}\norm{v}}{\abs{u^{\Htran}A^{\prime}(\lambda)v}}
        \quad\text{and}\quad 
        \kvec(A,\lambda)\defi \bignorm{(Q_{\perp}^{\Htran}A(\lambda)V_{\perp})^{-1}},
    \end{equation}
    where $Q_{\perp}$ and $V_{\perp}$ are orthonormal bases of $\spa{A^{\prime}(\lambda)v}^{\perp}$ and $\spa{v}^{\perp}$, respectively.

\end{lemma}

\begin{proof}
    By \cref{lem:exact}, we know that $(\lambda,e_{1})$ is a simple eigenpair of $B(\cdot)$ almost surely. We further denote its corresponding unit left eigenvector by $x$.
    
    \emph{Eigenvalue condition number:}
    Since $(0,v)$ is a simple eigenpair of the matrix $A(\lambda)$ and $[v,W_{\perp}]$ is orthonormal, we can let $P\in\C^{n\times (m-1)}$ be an orthonormal basis of $\range(A(\lambda)W_{\perp})$ and $p_{\perp} = (I-PP^{\Htran})A^{\prime}(\lambda)v$. Since $u^{\Htran}A(\lambda)=0$, we know that $u^{\Htran}P=0$ and 
    \begin{equation*}
        \norm{p_{\perp}} \geq \norm{uu^{\Htran}(I-PP^{\Htran})A^{\prime}(\lambda)v}= \norm{uu^{\Htran}A^{\prime}(\lambda)v} =  \abs{u^{\Htran}A^{\prime}(\lambda)v} = \frac{1}{\kval(A,\lambda)}.
    \end{equation*}
    Since the eigenvalue $0$ of the matrix $B(\lambda)$ is also simple, we know that  
    \begin{equation*}
        \spa{x} = \range\bigl(B(\lambda)\bigr)^{\perp}=\range\bigl(\Omega^{\Htran}A(\lambda)W_{\perp}\bigr)^{\perp}=\range(\Omega^{\Htran}P)^{\perp},    
    \end{equation*}
    and in turn, $x^{\Htran}\Omega^{\Htran} p_{\perp} = x^{\Htran}\Omega^{\Htran} A^{\prime}(\lambda)v-x^{\Htran}\Omega^{\Htran} PP^{\Htran}A^{\prime}(\lambda)v = x^{\Htran}\Omega^{\Htran} A^{\prime}(\lambda)v$.
    Note that $\Omega^{\Htran}P$ and $\Omega^{\Htran}p_{\perp}$ are independent; consequently, $x$ and $\Omega^{\Htran}p_{\perp}$ are also independent. Thus, 
    \begin{equation*}
        \frac{1}{\kval(B,\lambda)\cdot\norm{p_{\perp}}} = \frac{\abs{x^{\Htran}B^{\prime}(\lambda)e_{1}}}{\norm{p_{\perp}}} = \frac{\abs{x^{\Htran}\Omega^{\Htran}A^{\prime}(\lambda)v}}{\norm{p_{\perp}}} = \frac{\abs{x^{\Htran}\Omega^{\Htran}p_{\perp}}}{\norm{p_{\perp}}} = \Bigabs{x^{\Htran}\Omega^{\Htran}\bigl(p_{\perp}/\norm{p_{\perp}}\bigr)}
    \end{equation*}
    is the magnitude of a complex Gaussian random variable. Using the tail bound in \cref{lem:Gaussian}, we know that, with probability at least $1-\delta$, it holds that 
    \begin{equation*}
        \kval(B,\lambda)\leq  \frac{1}{\sqrt{\delta}\cdot\norm{p_{\perp}}}  \leq \frac{1}{\sqrt{\delta}}\cdot\kval(A,\lambda).
    \end{equation*}

    \emph{Eigenvector condition number:} Let 
    \begin{equation*}
        Q_{\perp}Q_{\perp}^{\Htran}A(\lambda)W_{\perp} = QR
    \end{equation*}
    be the thin QR decomposition. Since the matrix $[A^{\prime}(\lambda)v,A(\lambda)V_{\perp}]$ is nonsingular and $\range(W_{\perp})\subset \range(V_{\perp})$. We know that $Q\in\C^{n\times (m-1)}$ and
    \begin{equation*}
        \norm{R^{-1}} \leq \bignorm{\bigl(Q_{\perp}^{\Htran}A(\lambda)V_{\perp}\bigr)^{-1}} = \kvec(A,\lambda),
    \end{equation*}
    where we use the interlacing property on the smallest singular values of $Q_{\perp}^{\Htran}A(\lambda)W_{\perp}$ and $Q_{\perp}^{\Htran}A(\lambda)V_{\perp}$. 
    Recall that 
    \begin{equation*}
        \kvec(B,\lambda) =  \bignorm{\bigl((\Omega Z_{\perp})^{\Htran}A(\lambda)W_{\perp}\bigr)^{-1}}, 
    \end{equation*}
    where $Z_{\perp}$ is an orthonormal basis of $\spa{\Omega^{\Htran}A^{\prime}(\lambda)v}^{\perp}$. We know that $(\Omega Z_{\perp})^{\Htran}A^{\prime}(\lambda)v=0$, hence $\Omega Z_{\perp} = Q_{\perp}Q_{\perp}^{\Htran}\Omega Z_{\perp}$. Thus, 
    \begin{equation*}
        \begin{aligned}
            \kvec(B,\lambda) &= 
        \bignorm{\bigl((\Omega Z_{\perp})^{\Htran}Q_{\perp}Q_{\perp}^{\Htran}A(\lambda)W_{\perp}\bigr)^{-1}} = \bignorm{\bigl((\Omega Z_{\perp})^{\Htran}QR\bigr)^{-1}} \\ 
        &\leq \bignorm{\bigl(Z_{\perp}^{\Htran}(\Omega^{\Htran}Q)\bigr)^{-1}}\cdot \norm{R^{-1}}\leq \bignorm{\bigl(Z_{\perp}^{\Htran}(\Omega^{\Htran}Q)\bigr)^{-1}}\cdot \kvec(A,\lambda).        
        \end{aligned}
    \end{equation*}
    Note that $A^{\prime}(\lambda)v$ and $Q$ are orthogonal. We know that $\Omega^{\Htran}A^{\prime}(\lambda)v$ and $\Omega^{\Htran}Q$ are independent; consequently, $Z_{\perp}$ and $\Omega^{\Htran}Q$ are also independent. Then $Z_{\perp}^{\Htran}(\Omega^{\Htran}Q)\in\C^{(m-1)\times (m-1)}$ is a complex Gaussian random matrix. By \cref{lem:Gaussian},  
    \begin{equation*}
        \bignorm{\bigl(Z_{\perp}^{\Htran}(\Omega^{\Htran}Q)\bigr)^{-1}}\leq \sqrt{(m-1)/\delta}
    \end{equation*} 
    holds with probability at least $1-\delta$. 
\end{proof}


\subsection{Main result: convergence of randomized Ritz pairs}
\begin{theorem}
\label{thm:ritzpair}
Let $(\lambda,v)$ be a simple eigenpair of $A(\cdot)$.
For $0<\delta<1/4$, define
\begin{equation}
    \label{eq:defCr}
\Crval \defi 
\frac{C_{0}}{\sqrt{\delta}}
\cdot \kval(A,\lambda),
\quad\text{and}\quad 
\Crvec \defi  
1 + \frac{C_{0}\sqrt{m-1}}{\sqrt{\delta}}
\cdot \kvec(A,\lambda),
\end{equation}
where $C_{0}=\bigl(\sqrt{m}+\sqrt{2}+\sqrt{\ln(2/\delta)}\bigr)\norm{A(\lambda)}$, and $\kval(A,\lambda)$ and $\kvec(A,\lambda)$ are defined in~\cref{eq:defkappaA}.
Then there exist constants $\epsilon_{0}>0$ and $C>0$, depending only on $m$, $n$, $\delta$, $\kval(A,\lambda)$, $\kvec(A,\lambda)$, $\norm{A}_{\region}$, $\norm{A^{\prime}}_{\region}$, $\norm{A^{\prime\prime}}_{\region}$,
such that the following statement holds. 


Consider an arbitrary fixed orthonormal matrix $W_{\epsilon}\in\C^{n\times m}$ with $\epsilon=\angle(v,W_{\epsilon})\leq \epsilon_{0}$.
Let $B_{\epsilon}(\xi)=\Omega^{\Htran}A(\xi)W_{\epsilon}$, where $\Omega\in\C^{n\times m}$ is a complex Gaussian random matrix. Then, with probability at least $1-4\delta$, there exists an eigenpair $(\mu_{\epsilon},y_{\epsilon})$ of $B_{\epsilon}(\cdot)$ such that
\begin{equation}
    \label{eq:convRRR}
    \abs{\mu_{\epsilon}-\lambda} \leq \Crval\epsilon + C\epsilon^{2}
    \quad\text{and}\quad
    \angle(v,W_{\epsilon}y_{\epsilon}) \leq \Crvec\epsilon+ C\epsilon^{2}.
\end{equation}
\end{theorem}


\begin{remark}
    \label{rmk:nonasymptotic}
To achieve linear convergence, the coefficients $\{\Crval,\Crvec,C\}$ must be independent of the choice of $W_{\epsilon}$. Moreover, the failure event depends on $W_{\epsilon}$. Consequently, \cref{lem:pert}, on which \cref{thm:ritzpair} relies, cannot be replaced by asymptotic results, since we \textbf{cannot} take a union bound over a \textbf{family} of bases $W_{\epsilon}$.
\end{remark}


\begin{proof}[Proof of \cref{thm:ritzpair}]
    Since the choice of orthonormal bases of $\range(W_{\epsilon})$ does not affect randomized Ritz pairs, without loss of generality, we assume
    \begin{equation*}
        W_{\epsilon} = 
        \begin{bmatrix}
            v\cos\epsilon+w\sin\epsilon & W_{\perp}
        \end{bmatrix},
        \quad\text{where}\quad 
        \begin{bmatrix}
            v,w,W_{\perp}
        \end{bmatrix}
        \quad\text{is orthonormal.}
    \end{equation*}
    Let $B(\xi)=\Omega^{\Htran}A(\xi)W$, where $W=[v,W_{\perp}]$.
    Taking a union bound over \cref{lem:Gaussian,lem:exact,lem:kappa}, we know that, with probability at least $1-4\delta$, the following statements hold:
    \begin{equation}
        \label{eq:Fval}
        \begin{aligned}
            &\text{$B(\cdot)$ is regular and $(\lambda,e_{1})$ is its simple eigenpair},\\ 
            &\norm{\Omega}\leq \sqrt{n}+\sqrt{m}+\sqrt{\ln(2/\delta)},\\ 
            &\bignorm{\Omega^{\Htran}A(\lambda)[v,w]}\leq C_{0},\\
            & \kval(B,\lambda) \leq  \frac{1}{\sqrt{\delta}}\cdot \kval(A,\lambda),\\ 
            & \kvec(B,\lambda) \leq \frac{\sqrt{m-1}}{\sqrt{\delta}}\cdot \kvec(A,\lambda).
        \end{aligned}
    \end{equation}
    
    In what follows, we assume that \cref{eq:Fval} holds and apply \cref{lem:pert} to $B(\cdot)$ and $B_{\epsilon}(\cdot)$.
    By submultiplicativity, we know that, for any integer $k\geq 0$ and $\xi\in\region$,
    \begin{equation*}
        \Bigabs{\frac{\de^{k}}{\de \xi^{k}} e_{i}^{\Htran}B(\xi)e_{j}} = \Bigabs{(\Omega e_{i})^{\Htran} \frac{\de^{k}}{\de \xi^{k}} A(\xi ) (We_{j})} \leq \norm{\Omega} \bignorm{ \frac{\de^{k}}{\de \xi^{k}} A(\xi )},
    \end{equation*}
    implying that $B(\cdot)$ satisfies \cref{asp:A}. In particular, 
    \begin{equation*}
        \max\bigl\{\norm{B}_{\region},\norm{B^{\prime}}_{\region},\norm{B^{\prime\prime}}_{\region}\bigr\}\leq \bigl(\sqrt{n}+\sqrt{m}+\sqrt{\ln(2/\delta)}\bigr)\cdot 
        \max\bigl\{\norm{A}_{\region},\norm{A^{\prime}}_{\region},\norm{A^{\prime\prime}}_{\region}\bigr\}.
    \end{equation*}
    Let $\Delta B(\xi) \defi B_{\epsilon}(\xi)-B(\xi) =\Omega^{\Htran}A(\xi)(W_{\epsilon}-W)$. Note that
    \begin{equation}
        \label{eq:WeW}
        W_{\epsilon}-W  = \begin{bmatrix}
            v(\cos\epsilon-1)+w\sin\epsilon & 0_{n\times (m-1)}
        \end{bmatrix}
        = \begin{bmatrix}
            \begin{bmatrix}
                v & w
            \end{bmatrix}
            \begin{bmatrix}
                \cos\epsilon-1\\ 
                \sin\epsilon
            \end{bmatrix}& 0_{n\times (m-1)}
        \end{bmatrix}.
    \end{equation}
    Since  $(\cos\epsilon-1)^{2}+\sin^{2}\epsilon=4\sin^{2}(\epsilon/2)\leq \epsilon^{2}$, we know that $\norm{W_{\epsilon}-W}\leq \epsilon$, hence
    \begin{equation*}
        \max\bigl\{\norm{\Delta B}_{\region},\norm{\Delta B^{\prime}}_{\region},\norm{\Delta B^{\prime\prime}}_{\region}\bigr\}\leq \bigl(\sqrt{n}+\sqrt{m}+\sqrt{\ln(2/\delta)}\bigr)\cdot 
        \max\bigl\{\norm{A}_{\region},\norm{A^{\prime}}_{\region},\norm{A^{\prime\prime}}_{\region}\bigr\}\cdot \epsilon.
    \end{equation*}
    When $\epsilon_{0}$ is sufficiently small, \cref{lem:pert} yields the existence of an eigenpair $(\mu_{\epsilon},y_{\epsilon})$ of $B_{\epsilon}(\cdot)$ satisfying
    \begin{equation}
        \label{eq:pfpair}
        \begin{aligned}
            \abs{\mu_{\epsilon}-\lambda} &\leq \kval(B,\lambda)\cdot\norm{\Delta B(\lambda)} +C\epsilon^{2}\leq \frac{\Crval}{C_{0}}\norm{\Delta B(\lambda)} +C\epsilon^{2},\\ 
            \tan\angle(e_{1},y_{\epsilon})&\leq \kvec(B,\lambda)\cdot\norm{\Delta B(\lambda)} +C\epsilon^{2} \leq \frac{\Crvec-1}{C_{0}}\norm{\Delta B(\lambda)} +C\epsilon^{2},
        \end{aligned}
    \end{equation}
    where $C$ is a constant depending only on $m$, $n$, $\delta$, $\kval(A,\lambda)$, $\kvec(A,\lambda)$, $\norm{A}_{\region}$, $\norm{A^{\prime}}_{\region}$, $\norm{A^{\prime\prime}}_{\region}$.
    According \cref{eq:WeW}, it holds that 
    \begin{equation*}
        \norm{\Delta B(\lambda)}  
        =\Bignorm{\Omega^{\Htran}A(\lambda)\begin{bmatrix}
            v & w
        \end{bmatrix}
        \begin{bmatrix}
            \cos\epsilon-1\\ 
            \sin\epsilon
        \end{bmatrix}}\leq \Bignorm{\Omega^{\Htran}A(\lambda)\begin{bmatrix}
            v & w
        \end{bmatrix}} \cdot\epsilon \leq C_{0}\epsilon,
    \end{equation*}
    where we use \cref{eq:Fval} in the second inequality.
    Combining it with \cref{eq:pfpair}, we obtain the convergence of randomized Ritz values in \cref{eq:convRRR} and 
    \begin{equation*}
        \tan\angle(e_{1},y_{\epsilon})\leq (\Crvec-1)\epsilon +C\epsilon^{2}.
    \end{equation*}
    Then the convergence of randomized Ritz vectors in \cref{eq:convRRR} is proved by triangle inequality:
    \begin{equation}
        \angle(v,W_{\epsilon}y_{\epsilon})\leq 
        \angle(v,W_{\epsilon}e_{1})+\angle(W_{\epsilon}e_{1},W_{\epsilon}y_{\epsilon}) \leq \epsilon+\tan\angle(e_{1},y_{\epsilon}) = \Crvec\epsilon+C\epsilon^{2},
    \end{equation}
    where we use $\spa{W_{\epsilon}e_{1}}=\spa{W_{\epsilon}W_{\epsilon}^{\Htran}v}$ in the second inequality.
\end{proof}


\subsection{Refining randomized Ritz values}
\label{sec:refine}
This subsection investigates the accuracy of refined randomized Ritz values obtained from strategies described in \Cref{sec:algo}.
With the convergence of randomized Ritz vectors in \cref{thm:ritzpair}, the convergence of refined randomized Ritz values is actually straightforward. 
Before presenting results for general situations, we first state the result for standard eigenvalue problems.

\begin{theorem}
Let $(\lambda,v)$ be a simple eigenpair of a matrix $A_{0}\in\C^{n\times n}$. For $0<\delta<1/4$, we define $\Crvec$, $C$ and $\epsilon_{0}$ as in \cref{thm:ritzpair}. 
Consider an arbitrary fixed orthonormal matrix $W_{\epsilon}\in\C^{n\times m}$ with $\epsilon=\angle(v,W_{\epsilon})\leq \epsilon_{0}$.
Let $\Omega\in\C^{n\times m}$ be a complex Gaussian random matrix. With probability at least $1-4\delta$, there exists an eigenpair $(\mu_{\epsilon},y_{\epsilon})$ of the matrix pencil $\Omega^{\Htran}A_{0}W_{\epsilon}-\xi \Omega^{\Htran}W_{\epsilon}$ such that
\begin{equation*}
    \abs{\rho_{\epsilon} - \lambda}
    \leq \norm{A_{0} - \lambda I}\,\tan^{\gamma} \bigl(\Crvec \epsilon+C\epsilon^{2}\bigr),\quad \text{where}\quad \rho_{\epsilon} = \frac{(W_{\epsilon}y_{\epsilon})^{\Htran}A_{0}(W_{\epsilon}y_{\epsilon})}{(W_{\epsilon}y_{\epsilon})^{\Htran}(W_{\epsilon}y_{\epsilon})},
\end{equation*}
where $\gamma = 2$ if $A_{0}$ is Hermitian and $\gamma = 1$ otherwise.

\end{theorem}

The proof follows by substituting the convergence result for randomized Ritz vectors in \cref{thm:ritzpair} into the standard analysis of compression methods; see, for example, \cite[Chap.~4.3]{Saad2011}.  
Compared with randomized Ritz values $\mu_{\epsilon}$, the refinement $\rho_{\epsilon}$ leverages the Hermiticity of $A_{0}$ to achieve quadratic convergence, analogous to the standard Rayleigh--Ritz procedure for exterior eigenvalues.  
Moreover, when $v$ is (nearly) a left eigenvector of $A_{0}$ corresponding to $\lambda$, the accuracy of $\rho_{\epsilon}$ is substantially higher than that of $\mu_{\epsilon}$.





\paragraph{Refinement via Rayleigh functional.} 
\!We adopt the following lemma from \cite[Cor.~18 and Thm.~21]{Schwetlick2012}, which establishes both existence and accuracy of the Rayleigh functional for a general (not necessarily Hermitian) matrix-valued function.
\begin{lemma}
    Let $(\lambda,v)$ be a simple eigenpair of $A(\cdot)$. 
    If $v^{\Htran}A^{\prime}(\lambda)v\neq 0$, then there exists a unique Rayleigh functional $\Rq(\cdot)$ defined in a neighborhood of $v$. Moreover, for a vector $w_{\epsilon}$ with $\epsilon\defi\angle(v,w_{\epsilon})$, when $\epsilon$ is sufficiently small, it holds that 
    \begin{equation*}
        \abs{\Rq(w_{\epsilon}) - \lambda}
        \leq C_{\gamma}
        \frac{\norm{A(\lambda)}\norm{v}^{2}}{\abs{v^{\Htran}A^{\prime}(\lambda)v}}\tan^{\gamma}\epsilon,
    \end{equation*}
    where $\gamma = 2$ and $C_{\gamma} = 8/3$ if the matrix $A(\lambda)$ is Hermitian, while $\gamma = 1$ and $C_{\gamma} = 10/3$ otherwise.
\end{lemma}


Combining it with our convergence result of randomized Ritz vectors in \cref{thm:ritzpair}, we obtain the following convergence result on refined randomized Ritz values.


\begin{theorem}
    \label{thm:Rf}
    Let $(\lambda,v)$ be a simple eigenpair of $A(\cdot)$.
    Assume that $v^{\Htran}A^{\prime}(\lambda)v\neq 0$. Then there exists a unique Rayleigh functional $\Rq(\cdot)$ defined in a neighborhood of $v$.
    For $0<\delta<1/4$, we define $\Crvec$, $C$ and $\epsilon_{0}$ as in \cref{thm:ritzpair}. 
    Consider an arbitrary fixed orthonormal matrix $W_{\epsilon}\in\C^{n\times m}$ with $\epsilon=\angle(v,W_{\epsilon})\leq \epsilon_{0}$.
    Let $B_{\epsilon}(\xi)=\Omega^{\Htran}A(\xi)W_{\epsilon}$, where $\Omega\in\C^{n\times m}$ is a complex Gaussian random matrix. With probability at least $1-4\delta$, there exists an eigenpair $(\mu_{\epsilon},y_{\epsilon})$ of $B_{\epsilon}(\cdot)$ such that
    \begin{equation*}
        \abs{\Rq(W_{\epsilon}y_{\epsilon})-\lambda}\leq  
        C_{\gamma}\frac{\norm{A(\lambda)}\norm{v}^{2}}{\abs{v^{\Htran}A^{\prime}(\lambda)v}}
        \tan^{\gamma} \bigl(\Crvec \epsilon+C\epsilon^{2}\bigr),
    \end{equation*}
    where $\gamma = 2$ and $C_{\gamma} = 8/3$ if the matrix $A(\lambda)$ is Hermitian, while $\gamma = 1$ and $C_{\gamma} = 10/3$ otherwise.
\end{theorem}



\paragraph{Refinement via stationary points.}
As noted in \Cref{sec:algo}, a Rayleigh functional may not exist for certain matrix-valued functions. In such cases, one can instead refine randomized Ritz values by finding stationary points of \cref{eq:nls}. The following lemma demonstrates that the convergence of randomized Ritz vectors implies the convergence of the corresponding refinement.


\begin{lemma}
    \label{lem:ift}
    Let $(\lambda,v)$ be a simple eigenpair of $A(\cdot)$. Given a vector $w_{\epsilon}$ with $\epsilon\defi\angle(v,w_{\epsilon})$, when $\epsilon$ is sufficiently small, there exists a stationary point 
    \begin{equation*}
        \rho_{\epsilon}\in \Bigl\{ \rho\in\C \Bigm\vert \frac{\de }{\de \bar{\rho}} \norm{A(\rho)w_{\epsilon}}^{2} = 0  \Bigr\},
    \end{equation*}
    and a constant $C_{A}$ depending only on $A(\cdot)$, such that 
    \begin{equation*}
        \abs{\rho_{\epsilon}-\lambda} \leq  \frac{\norm{A(\lambda)}\norm{v}}{\norm{A^{\prime}(\lambda)v}}\epsilon+C_{A}\epsilon^{2}.
    \end{equation*}
\end{lemma}

The proof is based on the implicit function theorem and Taylor expansion; we defer it to \cref{app:ift}. By plugging the convergence of randomized Ritz vectors established in \cref{thm:ritzpair} into \cref{lem:ift}, we obtain the convergence of refined randomized Ritz values from finding stationary points of \cref{eq:nls}.



\begin{theorem}
    \label{thm:nls}
    Let $(\lambda,v)$ be a simple eigenpair of $A(\cdot)$.
    For $0<\delta<1/4$, we define $\Crvec$, $C$ and $\epsilon_{0}$ as in \cref{thm:ritzpair}. 
    Consider an arbitrary fixed orthonormal matrix $W_{\epsilon}\in\C^{n\times m}$ with $\epsilon=\angle(v,W_{\epsilon})\leq \epsilon_{0}$.
    Let $B_{\epsilon}(\xi)=\Omega^{\Htran}A(\xi)W_{\epsilon}$, where $\Omega\in\C^{n\times m}$ is a complex Gaussian random matrix. With probability at least $1-4\delta$, there exists an eigenpair $(\mu_{\epsilon},y_{\epsilon})$ of $B_{\epsilon}(\cdot)$ such that
    \begin{equation*}
        \abs{\rho_{\epsilon}-\lambda}\leq  
        \frac{\norm{A(\lambda)}\norm{v}}{\norm{A^{\prime}(\lambda)v}}
        \Crvec \epsilon+C_{A}(\Crvec^{2}+C)\epsilon^{2}
    \end{equation*}
    holds for some stationary point $\rho_{\epsilon}$ of \cref{eq:nls}, where $C_{A}$ is a constant depending only on $A(\cdot)$.
\end{theorem}







\section{Numerical experiments}
\label{sec:numexp}
In this section, we present several numerical experiments\footnote{All numerical experiments were implemented in MATLAB~R2022b and carried out on an AMD Ryzen~9 6900HX processor (8 cores, 3.3--4.9GHz) with 32GB of RAM. Scripts for reproducing the numerical results are publicly available at \url{https://github.com/nShao678/randomized-Rayleigh--Ritz}}.
Let $(\lambda,v)$ denote a target eigenpair. The accuracy of an approximate eigenpair $(\mu,w)$ is measured by
\begin{equation}
    \label{exp:err}
    \abs{\lambda-\mu} \quad\text{and}\quad \angle(v,w),
\end{equation}
respectively.
The experiments are designed to examine different aspects of the RRR procedure. First, we consider neutral modes arising in Hamiltonian systems, where the standard Rayleigh--Ritz procedure fails while the RRR procedure remains robust. Second, we solve the \texttt{butterfly} nonlinear eigenvalue problem from NLEVP \cite{Betcke2013} and illustrate that when $A(\lambda)$ is Hermitian, refined randomized Ritz values exhibit quadratic convergence comparable to that of (standard) Ritz values. Finally, we investigate the empirical failure probability of the RRR procedure, illustrating the sharpness of the factor $1/\sqrt{\delta}$ in \cref{thm:ritzpair}. 
This experiment also shows that this factor cannot be improved by oversampling.



\subsection{Neutral modes in Hamiltonian systems}
Consider a matrix pencil $A_{0}-\xi A_{1}\in\C^{2n\times 2n}$, where 
\begin{equation*}
    A_{0} = \begin{bmatrix}
        A_{11} & -A_{21}^{\Htran}\\ 
        A_{21} & A_{22}
    \end{bmatrix}
    \quad\text{and}\quad 
    A_{1}=\begin{bmatrix}
    0 & I\\ 
    I & 0
\end{bmatrix}
\end{equation*}
with $A_{11}$ and $A_{22}$ being Hermitian. Such generalized eigenvalue problems are often referred to as Hamiltonian eigenvalue problems \cite{Benner2005}. We are interested in the neutral modes \cite{Bognar1974}, that is, an eigenvector $v$ of the pencil $A_{0}-\xi A_{1}$ satisfying $v^{\Htran} A_{1} v = 0$.
Computing these neutral modes is important because they capture the directions in which the quadratic form associated with $A_{1}$ vanishes, highlighting marginally stable behavior in the system. Identifying these modes allows for a more accurate analysis of the system's critical dynamics and helps preserve the Hamiltonian structure in numerical computation, preventing spurious growth or decay of nearly neutral modes.

In this experiment, we consider a smoothly parameterized $A_{0}(\tau)$. Let $\lambda(\tau)$ be a branch of eigenvalues of $A_{0}(\tau)-\lambda A_{1}$ associated with neutral modes $v(\tau)$. Suppose that we have computed $v(\tau_{k})$ for some integers $k$, and our goal is to approximate $v(\tau_{0})$ and $\lambda(\tau_{0})$. In the following, we omit the parameter when it equals $\tau_{0}$. Let $\mathcal{W}_{k}=\spa{v(\tau_{1}),\dotsc,v(\tau_{k})}$. Then $\angle(v,\mathcal{W}_{k})$ is usually very small for a smooth $A_{0}(\cdot)$. However, since $v(\tau_{k})$ are all neutral modes, the standard Rayleigh--Ritz procedure fails in exact arithmetic as the compressed matrix pencil vanishes. In the following, we show that RRR works well for this task.

We construct an example as follows: Let 
\begin{equation*}
    Q = \frac{1}{2}\begin{bmatrix}
        Q_{1}+Q_{2} & Q_{1}-Q_{2}\\ 
        Q_{1}-Q_{2} & Q_{1}+Q_{2}
    \end{bmatrix}\in\C^{2n\times 2n},
\end{equation*}
where $Q_{1},Q_{2}\in\C^{n\times n}$ are Gaussian unitary ensembles.
Here, a Gaussian unitary ensemble is constructed from the $Q$-factor of the QR decomposition of a (square) complex Gaussian random matrix. 
Then $Q$ is also a unitary matrix and $Q^{\Htran}A_{1}Q=A_{1}$. 
Given a unit vector $v_{1}\in\C^{n}$, we assume the neutral modes admit 
\begin{equation*}
    v(\tau) = \eta(\tau)Q\begin{bmatrix}
        \exp\bigl((\tau-\tau_{0}) G\bigr)v_{1}\\ 0
    \end{bmatrix}
\end{equation*}
where $G\in\C^{n\times n}$ is a complex Gaussian random matrix.
Here $\eta(\tau)$ is a normalization factor such that $\norm{v(\tau)}=1$. Let 
\begin{equation}
    \label{eq:defHam}
    A_{0} = Q\begin{bmatrix}
        (I-v_{1}v_{1}^{\Htran})(G_{11}+G_{11}^{\Htran})(I-v_{1}v_{1}^{\Htran}) &-\overline{\lambda} v_{1}v_{1}^{\Htran}-(I-v_{1}v_{1}^{\Htran})G_{21}^{\Htran} \\ 
        \lambda v_{1}v_{1}^{\Htran}+G_{21}(I-v_{1}v_{1}^{\Htran}) &G_{22}+G_{22}^{\Htran}
    \end{bmatrix}Q^{\Htran},
\end{equation}
where $G_{11},G_{22}\in\C^{n\times n}$ are also complex Gaussian random matrices, $G_{21}$ will be specified later, and $\lambda$ is the target eigenvalue.
Then by construction, we know that $A_{0}v=\lambda A_{1}v$. 

Setting $n=2000$ and $\lambda=1$, we choose $v_{1}$ as a complex Gaussian random vector. We consider two choices of $G_{21}$, a zero matrix or a complex Gaussian random matrix. When $G_{21}=0$, then $v$ is both a left and a right eigenvector, while it is only a right eigenvector if $G_{21}$ is  complex Gaussian. Let $\tau_{k}=10^{-3}\cdot k$. We report numerical results for $k=1,\dotsc,10$ in \cref{fig:ham}. It is clear that (standard) Ritz values cannot provide meaningful approximation of the target eigenvalue. In exact arithmetic, (standard) Ritz vectors are not well-defined. Due to roundoff errors, the standard Rayleigh--Ritz procedure returns Ritz vectors, but they lose three digits of accuracy.
The RRR procedure works well in this situation: the accuracy of randomized Ritz vectors is nearly optimal, that is, they closely follow the subspace angles $\angle(v,\mathcal{W}_{k})$; and refined randomized Ritz values achieve quadratic convergence when $G_{21}=0$ and linear convergence for a complex Gaussian $G_{21}$.

\begin{figure}[htbp]
    \centering
    \includegraphics[width=\figsizeD]{ham2.\fig}
    \hspace{1cm}
    \includegraphics[width=\figsizeD]{ham1.\fig}
    \caption{Convergence of approximate neutral modes extracted from $\mathcal{W}_k$. The X-axis shows the dimension $k$ of $\mathcal{W}_{k}$, while the Y-axis reports the accuracy of approximate eigenpairs measured by \cref{exp:err}. Here, ``angle'' denotes $\angle(v,\mathcal{W}_k)$; ``RR-Vec/Val'' denote Ritz vectors/values from the standard Rayleigh--Ritz procedure; ``RRR-Vec/Val/reVal'' denote randomized Ritz vectors/values and refined randomized Ritz values from the RRR procedure (\cref{alg:RRR}). The matrix $G_{21}$ in \cref{eq:defHam} is zero in the left panel and a complex Gaussian random matrix in the right panel.} 
    \label{fig:ham}
\end{figure}


\subsection{The \texttt{butterfly} nonlinear eigenvalue problem}
Let  
\begin{equation*}
    A(\xi) = \sum_{i=0}^{4}c_{i}\xi^{i}A_{i}\in\R^{n\times n},
\end{equation*}
where $A_{0}$, $A_{2}$ and $A_{4}$ are symmetric and $A_{1}$ and $A_{3}$ are skew-symmetric.
Using the default parameters $c_{i}$ from \cite{Betcke2013} and $n=4096$, we compute the eigenvalue closest to $\sigma\in\{2\mi,1+1\mi\}$ with a nonlinear Arnoldi-type method \cite{Mehrmann2004}.
Specifically, we take $w_{0}$ as a complex Gaussian random starting vector, and then perform the residual inverse iteration \cite{Neumaier1985}: 
\begin{equation*}
    w_{k+1} = \eta_{k+1}\Bigl(w_{k}-\bigl(A(\sigma)\bigr)^{-1}A(\rho_{k})w_{k}\Bigr),
\end{equation*}
where $\eta_{k+1}$ is a normalization factor such that $\norm{w_{k+1}}=1$, and 
\begin{equation*}
    \rho_{k}=\Rq(w_{k}) = \argmin_{\rho\in\C} \bigl\{\abs{\rho-\sigma}\mid w_{k}^{\Htran}A(\rho)w_{k}=0 \bigr\}.
\end{equation*}
We construct a trial subspace as $\mathcal{W}_{k}=\spa{w_{1},\dotsc,w_{k}}$, and perform the standard and randomized Rayleigh--Ritz procedures to extract eigenpair approximations. The reference solution is computed by performing the residual inverse iteration until the residual norm is below $10^{-12}$.
Numerical results are reported in \cref{fig:butterfly}. 
When $\sigma=2\mi$, since the target eigenvalue $\lambda$ is purely imaginary, which makes $A(\lambda)$ Hermitian, both (standard) Ritz values and refined randomized Ritz values converge quadratically. In this situation, the behavior of the standard and randomized Rayleigh--Ritz procedures are similar. For the other situation, since $A(\lambda)$ is non-Hermitian, both methods exhibit linear convergence, and (refined) randomized Ritz values and vectors achieve approximately two additional digits of accuracy compared to (standard) Ritz values and vectors.


\begin{figure}[htbp]
    \centering
    \includegraphics[width=\figsizeD]{butterfly1.\fig}
    \hspace{1cm}
    \includegraphics[width=\figsizeD]{butterfly2.\fig}
    \caption{Convergence of approximate eigenpairs of the \texttt{butterfly} eigenvalue problem extracted from $\mathcal{W}_{k}$. The X-axis shows the dimension $k$ of $\mathcal{W}_{k}$, while the Y-axis reports the accuracy of approximate eigenpairs measured by \cref{exp:err}. Here, ``angle'' denotes $\angle(v,\mathcal{W}_k)$; ``RR-Vec/Val'' denote Ritz vectors/values from the standard Rayleigh--Ritz procedure; ``RRR-Vec/Val/reVal'' denote randomized Ritz vectors/values and refined randomized Ritz values from the RRR procedure (\cref{alg:RRR}). 
    The matrix $A(\lambda)$ is Hermitian for the left panel, and non-Hermitian for the right panel.}
    \label{fig:butterfly}
\end{figure}


\subsection{Empirical probability and oversampling}

\label{sec:numexpFP}

The prefactors $\Crval$ and $\Crvec$ in \cref{thm:ritzpair} both depend on $1/\sqrt{\delta}$. In certain randomized algorithms, such as the randomized SVD \cite{Halko2011}, this issue can be mitigated by oversampling; that is, by replacing the $n\times m$ complex Gaussian matrix $\Omega$ with an $n\times (m+s)$ complex Gaussian matrix $\Omega_{\mathrm{os}}$ for some integer $s \geq 1$.  
However, such a strategy is not effective in our setting. 

We illustrate that with a numerical experiment. Consider a matrix pencil $A_{0}-\xi A_{1}$, where $A_{0},A_{1}\in\C^{1000\times 1000}$ are both complex Gaussian random matrices. We are interested in the eigenvalue closest to $0.01$. Let $W_{\epsilon}\in\C^{1000\times 10}$ be an orthonormal basis computed by running $10$ steps of subspace shift-and-invert iterations from a complex Gaussian random starting block, where $\epsilon\approx 6.4\times 10^{-8}$. Then we extract approximate eigenpairs by the RRR procedure with and without oversampling. For oversampling, we use complex Gaussian matrices $\Omega_{\mathrm{os}}\in\C^{1000\times 20}$, that is $m=s=10$, and solve the rectangular eigenvalue problem of the pencil $\Omega_{\mathrm{os}}^{\Htran}A_{0}W_{\epsilon}-\xi \Omega_{\mathrm{os}}^{\Htran}W_{\epsilon}\in\C^{20\times 10}$ by the minimal perturbation approach proposed by \cite{Ito2016}, which essentially solves a total least squares problem via SVD. With $2^{17}$ independent trials on $\Omega$ and $\Omega_{\mathrm{os}}$, we collect the empirical result in \cref{fig:fp}, which indicates that the factor $1/\sqrt{\delta}$ is sharp and cannot be improved by oversampling. 

\begin{figure}[htbp]
    \centering
    \includegraphics[width=\figsizeD]{fp1.\fig}
    \hspace{1cm}
    \includegraphics[width=\figsizeD]{fp2.\fig}
    \caption{Empirical results from $2^{17}$ samples on accuracy (measured by \cref{exp:err}) for the RRR procedure (\cref{alg:RRR}).
    Here, ``RRR-Vec/Val/reVal'' denote randomized Ritz vectors/values and refined randomized Ritz values, and ``os'' denotes oversampling.
    The right panel reports the accuracy versus frequency.}
    \label{fig:fp}
\end{figure}




\section{Conclusion}





In this paper, we have introduced a randomized Rayleigh--Ritz (RRR) procedure for nonlinear eigenvalue problems.
Unlike the standard Rayleigh--Ritz approach, the RRR procedure reliably extracts approximate eigenpairs while achieving a convergence rate comparable to the ideal projection.
This addresses the long-standing eigenpair extraction problem posed in \cref{Q}.
Future work includes exploring ways to preserve Hermiticity for Hermitian problems and investigating how randomization can stabilize subspace methods in other contexts, such as evaluating matrix functions applied to vectors.

\section*{Acknowledgments}
The author thanks Daniel Kressner for helpful discussions and constructive feedback that significantly improved this work. 
\appendix

\section{Proof of \cref{lem:pert}}

\label{app:pfpert}
    Without loss of generality, we assume that $\norm{y}=1$.
    We first show that $Z_{\perp}^{\Htran}B(\lambda)Y_{\perp}$ is nonsingular via an argument similar to \cref{lem:exact}.
    Since $\lambda$ is a simple eigenvalue associated with left and right eigenvectors $x$ and $y$, we know that 
    \begin{equation*}
        \rk(B(\lambda)Y_{\perp}) = m-1,\quad x^{\Htran}B(\lambda)Y_{\perp}=0
        \quad\text{and}\quad x^{\Htran}B^{\prime}(\lambda)y\neq 0.
    \end{equation*}
    Thus, $[B^{\prime}(\lambda)y,B(\lambda)Y_{\perp}]\in\C^{m\times m}$ is nonsingular. In turn, $Z_{\perp}^{\Htran}B(\lambda)Y_{\perp}$ is nonsingular, where we recall that $Z_{\perp}$ is an orthonormal basis of $\spa{B^{\prime}(\lambda)y}^{\perp}$.

    Then we consider the analytic expansion.
    Since $\lambda$ is simple, \cite[Thm.~2.1]{Andrew1993} ensures that the eigenvalue $\lambda$ and its corresponding eigenvector $y$ are (locally) analytic with respect to the perturbation. 
    Thus, when $\epsilon$ is sufficiently small, we can assume $\lambda+t\Delta \lambda\in\region$ for all $0\leq t\leq 1$.
    Consider the vector-valued function 
    \begin{equation*}
        f(t) = B(\lambda+t\Delta \lambda)\cdot (y+t\Delta y) + t\Delta B(\lambda+t\Delta \lambda)\cdot (y+t\Delta y).
    \end{equation*}
    Since $f(0)=f(1)=0$, we know that $\int_{0}^{1}f^{\prime}(t)\de t=0$. Hence, 
    \begin{equation}
        \label{eq:appint}
        \norm{f^{\prime}(0)} = \Bignorm{\int_{0}^{1}f^{\prime}(0)-f^{\prime}(t) \de t}\leq \int_{0}^{1}\norm{f^{\prime}(0)-f^{\prime}(t)} \de t\leq \sup_{0\leq t\leq 1}\norm{f^{\prime\prime}(t)}. 
    \end{equation}
    By the definition of $f$, we have 
    \begin{equation}
        \label{eq:appTaylor}
        f^{\prime}(0) = \Delta\lambda \cdot B^{\prime}(\lambda) y+B(\lambda)\Delta y+\Delta B(\lambda)\cdot y.
    \end{equation}
    For the second order derivative, direct computation yields that 
    \begin{equation*}
        \begin{aligned}
            f^{\prime\prime}(t)  =& (\Delta \lambda)^{2}\cdot \bigl(B^{\prime\prime}(\lambda+t\Delta\lambda)+t\Delta B^{\prime\prime}(\lambda+t\Delta\lambda)\bigr)\cdot(y+t\Delta y)\\
            &+2(\Delta \lambda)\cdot \bigl(B^{\prime}(\lambda+t\Delta\lambda)+t\Delta B^{\prime}(\lambda+t\Delta\lambda)\bigr)\cdot\Delta y\\
            &+2(\Delta \lambda)\cdot \Delta B^{\prime}(\lambda+t\Delta\lambda)\cdot(y+t\Delta y)
            +2\Delta B(\lambda+t\Delta\lambda)\cdot\Delta y.
        \end{aligned} 
    \end{equation*}
    Applying the triangle inequality and \cref{eq:appint}, we have 
    \begin{equation}
        \label{eq:appPertPf1}
        \begin{aligned}
            \norm{f^{\prime}(0)}\leq \sup_{0\leq t\leq 1}\norm{f^{\prime\prime}(t)}\leq &(M+\epsilon)\abs{\Delta \lambda}^{2}(1+\norm{\Delta y})+2(M+\epsilon)\abs{\Delta \lambda}\norm{\Delta y}\\
            &+2\epsilon\abs{\Delta \lambda}(1+\norm{\Delta y})+2\epsilon\norm{\Delta y}.
        \end{aligned}
    \end{equation}
    Multiplying $x^{\Htran}$ from the left-hand side of \cref{eq:appTaylor} yields that 
    \begin{equation}
        \label{eq:appPertPf2}
        \abs{\Delta\lambda} = \frac{\abs{-x^{\Htran}\Delta B(\lambda)y+x^{\Htran}f^{\prime}(0)}}{\abs{x^{\Htran}B^{\prime}(\lambda)y}}\leq \kval(B,\lambda)\cdot\bigl(\norm{\Delta B(\lambda)} + \norm{f^{\prime}(0)}\bigr),
    \end{equation}
    where $x^{\Htran}B(\lambda)=0$ is used. Multiplying $Z_{\perp}^{\Htran}$ from the left-hand side of \cref{eq:appTaylor} yields
    \begin{equation*}
        Z_{\perp}^{\Htran}f^{\prime}(0) = 
        \Delta \lambda \cdot Z_{\perp}^{\Htran}B^{\prime}(\lambda)y+
        Z_{\perp}^{\Htran}B(\lambda)\begin{bmatrix}
            y&Y_{\perp}
        \end{bmatrix}
        \begin{bmatrix}
            y^{\Htran}\\ Y_{\perp}^{\Htran}
        \end{bmatrix}
        \Delta y  + Z_{\perp}^{\Htran}\Delta B(\lambda)y. 
    \end{equation*}
    Using $Z_{\perp}^{\Htran}B^{\prime}(\lambda)y=0$ and $B(\lambda)y=0$, we have 
    \begin{equation*}
        Z_{\perp}^{\Htran}B(\lambda)Y_{\perp}\cdot Y_{\perp}^{\Htran}\Delta y
         = Z_{\perp}^{\Htran}f^{\prime}(0)- Z_{\perp}^{\Htran}\Delta B(\lambda)y. 
    \end{equation*}
    Multiplying $(Z_{\perp}^{\Htran}B(\lambda)Y_{\perp})^{-1}$ from the left-hand side and using $\Delta y = Y_{\perp}Y_{\perp}^{\Htran}\Delta y$, we have 
    \begin{equation}
        \label{eq:appPertPf3}
        \norm{\Delta y} = \bignorm{(Z_{\perp}^{\Htran}B(\lambda)Y_{\perp})^{-1}Z_{\perp}^{\Htran}\bigl(\Delta B(\lambda)y-f^{\prime}(0)\bigr)} \leq \kvec(B,\lambda)\cdot \bigl(\norm{\Delta B(\lambda)}+\norm{f^{\prime}(0)}\bigr).
    \end{equation}
        
We complete the proof via a bootstrap argument. Write $\kappa_{1}=\kval(B,\lambda)$ and $\kappa_{2}=\kvec(B,\lambda)$, then we have  $M\kappa_{1}\geq 1$ and $M\kappa_{2}\geq 1$ by $\kappa_{1}\geq 1/\norm{B^{\prime}(\lambda)}$ and $\kappa_{2}\geq 1/\norm{B(\lambda)}$. In the following, we assume that 
\begin{equation}
    \label{eq:appPertEps}
    56M\kappa_{1}(\kappa_{1}+\kappa_{2})\epsilon\leq 1,
\end{equation}
which implies $\epsilon\leq M$ and $2\kappa_{2}\epsilon\leq 1$.

Replacing $\Delta B$ by $s\Delta B$ for $s\in[0,1]$ scales $\epsilon$ and $\norm{\Delta B(\lambda)}$ to $s\epsilon$ and $s\norm{\Delta B(\lambda)}$ in all estimates above, and \cite[Thm.~2.1]{Andrew1993} provides an analytic branch of eigenpairs $\bigl(\lambda+\Delta \lambda(s),y+\Delta y(s)\bigr)$ of $B+s\Delta B$ starting at $(\lambda,y)$. We use $f_{s}$ to denote the function $f$ formed with $s\Delta B$, $\Delta \lambda(s)$ and $\Delta y(s)$. Suppose that
\begin{equation}
    \label{eq:appPertPf4}
    \abs{\Delta \lambda(s)} \leq 2\kappa_{1}s\epsilon
    \quad\text{and}\quad
    \norm{\Delta y(s)} \leq 2\kappa_{2}s\epsilon.
\end{equation}
Then $\norm{\Delta y(s)}\leq 1$ and $s\epsilon\leq M$, so \cref{eq:appPertPf1} yields
\begin{equation}
    \label{eq:appPertPf5}
    \norm{f_{s}^{\prime}(0)} \leq \bigl(16M\kappa_{1}^{2}+16M\kappa_{1}\kappa_{2}+8\kappa_{1}+4\kappa_{2}\bigr)(s\epsilon)^{2}
    \leq 28M\kappa_{1}(\kappa_{1}+\kappa_{2})(s\epsilon)^{2},
\end{equation}
and substituting this into \cref{eq:appPertPf2,eq:appPertPf3} improves \cref{eq:appPertPf4} to
\begin{equation*}
    \abs{\Delta \lambda(s)}\leq \kappa_{1}s\epsilon\bigl(1+28M\kappa_{1}(\kappa_{1}+\kappa_{2})\epsilon\bigr)\leq \frac{3}{2}\kappa_{1}s\epsilon
    \quad\text{and}\quad
    \norm{\Delta y(s)}\leq \frac{3}{2}\kappa_{2}s\epsilon.
\end{equation*}
Since \cref{eq:appPertPf4} holds for small $s$ and the improvement above is strict, a continuity argument shows that \cref{eq:appPertPf4} holds for all $s\in[0,1]$. Taking $s=1$ in \cref{eq:appPertPf5} and substituting it back into \cref{eq:appPertPf2,eq:appPertPf3} completes the proof, with $C=28M\kappa_{1}(\kappa_{1}+\kappa_{2})\max\{\kappa_{1},\kappa_{2}\}$.
    
    

\section{Proof of \cref{lem:ift}}
\label{app:ift}

To avoid the proliferation of constants, we use $C_{A}$ to denote generic constants depending only on $A(\cdot)$, whose values may change from line to line. 

Let us first show that, as $\epsilon\to 0$, there exists a stationary point
\begin{equation}
    \label{eq:apprhoeps}
    \rho_{\epsilon} \in \Bigl\{ \rho\in\C \Bigm\vert \frac{\de }{\de \bar{\rho}} \norm{A(\rho)w_{\epsilon}}^{2} = 0  \Bigr\}, 
    \quad \text{such that}\quad 
    \abs{\rho_{\epsilon}-\lambda}\leq C_{A}\epsilon.
\end{equation}
Without loss of generality, we assume that $\norm{v}=\norm{w_{\epsilon}}=1$ and $\norm{v-w_{\epsilon}}=2\sin\epsilon/2\leq \epsilon$. Denote 
\begin{equation*}
    \lambda = \alpha_{0}+\beta_{0}\mi
    \quad\text{and}\quad 
    v = x_{0}+y_{0}\mi,
    \quad\text{where}\quad
    \alpha_{0},\beta_{0}\in\R
    \quad\text{and}\quad 
    x_{0},y_{0}\in\R^{n}.   
\end{equation*}
Consider the \emph{real} analytic function
\begin{equation*}
    f(\alpha,\beta,x,y) = \norm{A(\rho)w}^{2},
    \quad\text{where}\quad
    \rho=\alpha+\beta\mi 
    \quad\text{and}\quad 
    w = x+y\mi.
\end{equation*}
Since $A(\lambda)v=0$, we have $f(\alpha_{0},\beta_{0},x_{0},y_{0})=0$. Define the function $G\colon \R^{2n+2}\mapsto \R^{2}$ as 
\begin{equation*}
    G(\alpha,\beta,x,y) = \begin{bmatrix}
        \nabla_{\alpha} f(\alpha,\beta,x,y)\\ 
        \nabla_{\beta} f(\alpha,\beta,x,y)
    \end{bmatrix}.
\end{equation*}
We know $G(\alpha_{0},\beta_{0},x_{0},y_{0})=0$. Since $A(\cdot)$ is holomorphic, Taylor expansion yields that 
\begin{equation*}
    A(\lambda+\Delta \lambda) v = A(\lambda)v+\Delta \lambda\cdot A^{\prime}(\lambda) v +\order(\abs{\Delta \lambda}^{2}) = \Delta \lambda\cdot A^{\prime}(\lambda) v +\order(\abs{\Delta \lambda}^{2}).
\end{equation*}
Since $\lambda$ is a simple eigenvalue, we know $u^{\Htran}A^{\prime}(\lambda)v\neq 0$, implying that $A^{\prime}(\lambda)v\neq 0$. Thus, 
\begin{equation*}
    \norm{A(\lambda+\Delta \lambda)v}^{2} = \abs{\Delta \lambda}^{2}\norm{A^{\prime}(\lambda)v}^{2}+\order(\abs{\Delta \lambda}^{3}),
\end{equation*}
implying that the Jacobian matrix of $G$ is 
\begin{equation*}
    \nabla_{\alpha,\beta}^{2}f(\alpha_{0},\beta_{0},x_{0},y_{0}) = 2\norm{A^{\prime}(\lambda)v}^{2}I_{2}\succ 0.
\end{equation*}
The implicit function theorem asserts that there exist real analytical functions $f_{\alpha}(x,y)$ and $f_{\beta}(x,y)$ defined in a neighborhood of $(x_{0},y_{0})$ such that 
\begin{equation*}
    \alpha_{0}=f_{\alpha}(x_{0},y_{0}),\quad 
    \beta_{0}=f_{\beta}(x_{0},y_{0}),\quad 
    G\bigl(f_{\alpha}(x,y),f_{\beta}(x,y),x,y\bigr) = 0.
\end{equation*}
Here both $f_{\alpha}$ and $f_{\beta}$ are independent of $w_{\epsilon}$.
Since $\norm{w_{\epsilon}-v}\leq \epsilon$, we know that $\norm{\Re w_{\epsilon}-x_{0}}\leq \epsilon$ and $\norm{\Im w_{\epsilon}-y_{0}}\leq \epsilon$. When $\epsilon$ is sufficiently small, we can take 
\begin{equation*}
    \rho_{\epsilon} = f_{\alpha}(\Re w_{\epsilon},\Im w_{\epsilon})+f_{\beta}(\Re w_{\epsilon},\Im w_{\epsilon})\mi \in \Bigl\{ \rho\in\C \Bigm\vert \frac{\de }{\de \bar{\rho}} \norm{A(\rho)w_{\epsilon}}^{2} = 0  \Bigr\},
\end{equation*}
such that $\abs{\rho_{\epsilon}-\lambda}\leq C_{A}\epsilon$. Thus, \cref{eq:apprhoeps} is justified.

Then we further investigate $\abs{\rho_{\epsilon}-\lambda}$. By the stationary condition of \cref{eq:apprhoeps}, we have 
\begin{equation}
    \label{eq:apppfT1}
    \bigl(A^{\prime}(\rho_{\epsilon})w_{\epsilon}\bigr)^{\Htran}\bigl(A(\rho_{\epsilon})w_{\epsilon}\bigr) = 0.
\end{equation}
Taylor expansions of $A(\rho_{\epsilon})w_{\epsilon}$ and $A^{\prime}(\rho_{\epsilon})w_{\epsilon}$ give that
\begin{equation*}
    A(\rho_{\epsilon})w_{\epsilon} = (\rho_{\epsilon}-\lambda)A^{\prime}(\lambda)v+A(\lambda)(w_{\epsilon}-v)+r_{2}, 
    \quad\text{and}\quad 
    A^{\prime}(\rho_{\epsilon})w_{\epsilon} = A^{\prime}(\lambda)v +r_{1},        
\end{equation*}
where $\norm{r_{2}}\leq C_{A}\epsilon^{2}$ and $\norm{r_{1}}\leq C_{A}\epsilon$. 
Here, we use $\norm{v-w_{\epsilon}}\leq \epsilon$, $\abs{\lambda-\rho_{\epsilon}}\leq C_{A}\epsilon$ and $A(\lambda)v=0$. 
Substituting these two expansions into \cref{eq:apppfT1}, we have 
\begin{equation*}
    \Bigabs{
    (\rho_{\epsilon}-\lambda)\norm{A^{\prime}(\lambda)v}^{2}+\bigl(A^{\prime}(\lambda)v\bigr)^{\Htran}\bigl(A(\lambda)(w_{\epsilon}-v)\bigr)} \leq C_{A}\epsilon^{2}.
\end{equation*}
Rearranging this equation and using the submultiplicativity, we complete the proof:
\begin{equation*}
    \abs{\rho_{\epsilon}-\lambda}\leq \frac{\norm{A(\lambda)}}{\norm{A^{\prime}(\lambda)v}}\epsilon+C_{A}\epsilon^{2}.
\end{equation*}



\begin{thebibliography}{10}

\bibitem{Andrew1993}
A.~L. Andrew, K.-W. Chu, and P.~Lancaster.
\newblock Derivatives of eigenvalues and eigenvectors of matrix functions.
\newblock {\em SIAM J. Matrix Anal. Appl.}, 14(4):903--926, 1993.
\newblock \href {https://doi.org/10.1137/0614061} {\path{doi:10.1137/0614061}}.

\bibitem{Benner2005}
P.~Benner, D.~Kressner, and V.~Mehrmann.
\newblock Skew-{H}amiltonian and {H}amiltonian eigenvalue problems: theory, algorithms and applications.
\newblock In {\em Proceedings of the {C}onference on {A}pplied {M}athematics and {S}cientific {C}omputing}, pages 3--39. Springer, Dordrecht, 2005.
\newblock \href {https://doi.org/10.1007/1-4020-3197-1\_1} {\path{doi:10.1007/1-4020-3197-1\_1}}.

\bibitem{Betcke2013}
T.~Betcke, N.~J. Higham, V.~Mehrmann, C.~Schr\"oder, and F.~Tisseur.
\newblock N{LEVP}: a collection of nonlinear eigenvalue problems.
\newblock {\em ACM Trans. Math. Software}, 39(2):Art. 7, 28, 2013.
\newblock \href {https://doi.org/10.1145/2427023.2427024} {\path{doi:10.1145/2427023.2427024}}.

\bibitem{Bognar1974}
J.~Bogn\'ar.
\newblock {\em Indefinite inner product spaces}, volume Band 78 of {\em Ergebnisse der Mathematik und ihrer Grenzgebiete}.
\newblock Springer-Verlag, New York-Heidelberg, 1974.

\bibitem{Dedieu2000}
J.-P. Dedieu and F.~Tisseur.
\newblock Perturbation theory for homogeneous polynomial eigenvalue problems.
\newblock {\em Linear Algebra Appl.}, 358:71--94, 2003.
\newblock \href {https://doi.org/10.1016/S0024-3795(01)00423-2} {\path{doi:10.1016/S0024-3795(01)00423-2}}.

\bibitem{Edelman1988}
A.~Edelman.
\newblock Eigenvalues and condition numbers of random matrices.
\newblock {\em SIAM J. Matrix Anal. Appl.}, 9(4):543--560, 1988.
\newblock \href {https://doi.org/10.1137/0609045} {\path{doi:10.1137/0609045}}.

\bibitem{Guttel2017}
S.~G\"uttel and F.~Tisseur.
\newblock The nonlinear eigenvalue problem.
\newblock {\em Acta Numer.}, 26:1--94, 2017.
\newblock \href {https://doi.org/10.1017/S0962492917000034} {\path{doi:10.1017/S0962492917000034}}.

\bibitem{Halko2011}
N.~Halko, P.~G. Martinsson, and J.~A. Tropp.
\newblock Finding structure with randomness: probabilistic algorithms for constructing approximate matrix decompositions.
\newblock {\em SIAM Rev.}, 53(2):217--288, 2011.
\newblock \href {https://doi.org/10.1137/090771806} {\path{doi:10.1137/090771806}}.

\bibitem{Ito2016}
S.~Ito and K.~Murota.
\newblock An algorithm for the generalized eigenvalue problem for nonsquare matrix pencils by minimal perturbation approach.
\newblock {\em SIAM J. Matrix Anal. Appl.}, 37(1):409--419, 2016.
\newblock \href {https://doi.org/10.1137/14099231X} {\path{doi:10.1137/14099231X}}.

\bibitem{Jia1997}
Z.~Jia.
\newblock Refined iterative algorithms based on {A}rnoldi's process for large unsymmetric eigenproblems.
\newblock {\em Linear Algebra Appl.}, 259:1--23, 1997.
\newblock \href {https://doi.org/10.1016/S0024-3795(96)00238-8} {\path{doi:10.1016/S0024-3795(96)00238-8}}.

\bibitem{Jia2001}
Z.~Jia and G.~W. Stewart.
\newblock An analysis of the {R}ayleigh--{R}itz method for approximating eigenspaces.
\newblock {\em Math. Comp.}, 70(234):637--647, 2001.
\newblock \href {https://doi.org/10.1090/S0025-5718-00-01208-4} {\path{doi:10.1090/S0025-5718-00-01208-4}}.

\bibitem{Jia2025}
Z.~Jia and Q.~Zheng.
\newblock An analysis of the {R}ayleigh--{R}itz and refined {R}ayleigh-{R}itz methods for regular nonlinear eigenvalue problems.
\newblock {\em SIAM J. Matrix Anal. Appl.}, 46(1):676--701, 2025.
\newblock \href {https://doi.org/10.1137/23M161392X} {\path{doi:10.1137/23M161392X}}.

\bibitem{Mehrmann2004}
V.~Mehrmann and H.~Voss.
\newblock Nonlinear eigenvalue problems: a challenge for modern eigenvalue methods.
\newblock {\em GAMM Mitt. Ges. Angew. Math. Mech.}, 27(2):121--152, 2004.
\newblock \href {https://doi.org/10.1002/gamm.201490007} {\path{doi:10.1002/gamm.201490007}}.

\bibitem{Morgan1991}
R.~B. Morgan.
\newblock Computing interior eigenvalues of large matrices.
\newblock {\em Linear Algebra Appl.}, 154/156:289--309, 1991.
\newblock \href {https://doi.org/10.1016/0024-3795(91)90381-6} {\path{doi:10.1016/0024-3795(91)90381-6}}.

\bibitem{Nakatsukasa2024}
Y.~Nakatsukasa and J.~A. Tropp.
\newblock Fast and accurate randomized algorithms for linear systems and eigenvalue problems.
\newblock {\em SIAM J. Matrix Anal. Appl.}, 45(2):1183--1214, 2024.
\newblock \href {https://doi.org/10.1137/23M1565413} {\path{doi:10.1137/23M1565413}}.

\bibitem{Neumaier1985}
A.~Neumaier.
\newblock Residual inverse iteration for the nonlinear eigenvalue problem.
\newblock {\em SIAM J. Numer. Anal.}, 22(5):914--923, 1985.
\newblock \href {https://doi.org/10.1137/0722055} {\path{doi:10.1137/0722055}}.

\bibitem{Nocedal2006}
J.~Nocedal and S.~J. Wright.
\newblock {\em Numerical optimization}.
\newblock Springer Series in Operations Research and Financial Engineering. Springer, New York, second edition, 2006.

\bibitem{Rudelson2010}
M.~Rudelson and R.~Vershynin.
\newblock Non-asymptotic theory of random matrices: extreme singular values.
\newblock In {\em Proceedings of the {I}nternational {C}ongress of {M}athematicians. {V}olume {III}}, pages 1576--1602. Hindustan Book Agency, New Delhi, 2010.

\bibitem{Saad2011}
Y.~Saad.
\newblock {\em Numerical methods for large eigenvalue problems}, volume~66 of {\em Classics in Applied Mathematics}.
\newblock Society for Industrial and Applied Mathematics (SIAM), Philadelphia, PA, revised edition, 2011.
\newblock \href {https://doi.org/10.1137/1.9781611970739} {\path{doi:10.1137/1.9781611970739}}.

\bibitem{Schwetlick2012}
H.~Schwetlick and K.~Schreiber.
\newblock Nonlinear {R}ayleigh functionals.
\newblock {\em Linear Algebra Appl.}, 436(10):3991--4016, 2012.
\newblock \href {https://doi.org/10.1016/j.laa.2010.06.048} {\path{doi:10.1016/j.laa.2010.06.048}}.

\bibitem{Stewart2001}
G.~W. Stewart.
\newblock {\em Matrix algorithms. {V}ol. {II}}.
\newblock Society for Industrial and Applied Mathematics (SIAM), Philadelphia, PA, 2001.
\newblock \href {https://doi.org/10.1137/1.9780898718058} {\path{doi:10.1137/1.9780898718058}}.

\bibitem{Stewart1990}
G.~W. Stewart and J.~G. Sun.
\newblock {\em Matrix perturbation theory}.
\newblock Computer Science and Scientific Computing. Academic Press, Inc., Boston, MA, 1990.

\end{thebibliography}

\bibliographystyle{abbrvurl}
\end{document}
